\chapter{Functions}
\section{Preliminary Remarks}
The concept of a function is prevalent in mathematics. As such, there is a compelling need
to provide a rigorous definition for it. Fortunately, as we shall soon see,
it is entirely possible (and indeed sufficient) to represent functions as simply sets with
certain properties. Before a set-theoretic definition for functions can be stated, however,
a more fundamental concept has to be introduced:
\keyword{ordered pairs}.

\section{Ordered Pairs}

\begin{dfn}[Kuratowski's Definition]
    \label{dfn:kuratowski-ordered-pairs}
    For any two objects \( a \) and \( b \), the \keyword{ordered pair} \orderedpair{a}{b}
    is defined as the set \kuratowski{a}{b}.
    \begin{equation*}
        \orderedpair{a}{b} \coloneq \kuratowski{a}{b}
    \end{equation*}
    Using \cref{dfn:unordered-pair}, we can formulate this definition as,
    \begin{equation*}
        \forall a,b \, \forall x \, \mleft( x \in \orderedpair{a}{b}
        \iff
        x = \mleft\{ a \mright\}
        \lor
        x = \mleft\{ a, b \mright\}
        \mright)
    \end{equation*}
    Note that if \( a = b\), then \orderedpair{a}{a} = \kuratowski{a}{a} = \( \mleft\{ \mleft\{ a \mright\} \mright\} \).
    For any ordered pair \orderedpair{a}{b}, we refer to \( a \) and \( b \) as its first and second coordinates respectively.
\end{dfn}

\begin{thm}[Existence]
    \label{thm:existence-of-ordered-pair}
    An ordered pair exists.
\end{thm}

\begin{proof}
    By the \nameref{axm:axiom-of-existence}, the empty set \( \emptyset \) exists. Since \( \emptyset \) exists,
    then the \nameref{axm:axiom-of-pairing} guarantees that the set \( \mleft\{ \emptyset \mright\} \) exists.
    And since \( \mleft\{ \emptyset \mright\} \) exists, we apply the \nameref{axm:axiom-of-pairing} once again
    to infer that the set \( \mleft\{ \mleft\{ \emptyset \mright\} \mright\} \) exists. But by \cref{dfn:kuratowski-ordered-pairs},
    \begin{equation*}
        \orderedpair{\emptyset}{\emptyset} = \kuratowski{\emptyset}{\emptyset} = \mleft\{ \mleft\{ \emptyset \mright\} \mright\}
    \end{equation*}
    Thus the ordered pair \orderedpair{\emptyset}{\emptyset} exists.
\end{proof}

\begin{thm}[Ordered Pair Equality]
    \label{thm:equality-of-ordered-pairs}
    For any two ordered pairs \orderedpair{a}{b} and \orderedpair{c}{d}, \( \orderedpair{a}{b} = \orderedpair{c}{d} \)
    if and only if \( a = c \) and \( b = d \).
    \begin{equation*}
        \forall a,b,c,d \, \big( \orderedpair{a}{b} = \orderedpair{c}{d}
        \iff
        \mleft( a = c \mright) \land \mleft( b = d \mright) \big)
    \end{equation*}
\end{thm}

\begin{proof}
    Let \( \alpha = \orderedpair{a}{b} \) and \( \beta = \orderedpair{c}{d} \) be arbitrary ordered pairs.
    Thus, by \cref{dfn:kuratowski-ordered-pairs}, the following statements must hold:
    \begin{gather}
        \forall x \, \mleft( x \in \alpha \iff x = \mleft\{ a \mright\} \lor x = \mleft\{ a, b \mright\} \mright)
        \label{eq:predicate-for-alpha}\\
        \forall x \, \mleft( x \in \beta \iff x = \mleft\{ c \mright\} \lor x = \mleft\{ c, d \mright\} \mright) \label{eq:predicate-for-beta}
    \end{gather}
    Now, suppose that \( \alpha = \beta \). By the \nameref{axm:axiom-of-extensionality}, we have,
    \begin{equation}
        \label{eq:alpha-beta-biconditional}
        \forall x \, \mleft( x \in \alpha \iff x \in \beta \mright)\\
    \end{equation}
    By the \aref{rule-of-inference:law-of-excluded-middle}{Law of Excluded Middle}, obviously,
    \begin{equation*}
        a = b \lor a \neq b
    \end{equation*}
    {\scshape Case 1}: \( a = b \). From \eqref{eq:predicate-for-alpha},
    \( x \in \alpha \iff x = \mleft\{ a \mright\} \) for any \( x \).
    Clearly, \( \mleft\{ a \mright\} \in \alpha \),
    so by \eqref{eq:alpha-beta-biconditional}, we infer \( \mleft\{ a \mright\} \in \beta \).
    By \eqref{eq:predicate-for-beta}, we have \( \mleft\{ a \mright\} = \mleft\{ c \mright\} \lor \mleft\{ a \mright\} = \mleft\{ c, d \mright\} \).
    In either case, \( a = c \) is implied. Since \( a = b = c = d\), \( b = d \) also follows.

    \vspace{0.5cm}
    \noindent
    {\scshape Case 2}: \( a \neq b \). It follows from \eqref{eq:predicate-for-alpha} that \( \mleft\{ a \mright\} \in \alpha \)
    and \( \mleft\{ a, b \mright\} \in \alpha \). By \eqref{eq:alpha-beta-biconditional}, we have \( \mleft\{ a \mright\} \in \beta \)
    and \( \mleft\{ a, b \mright\} \in \beta \). Thus, the following statements are both true by \eqref{eq:predicate-for-beta},
    \begin{gather}
        \mleft\{ a \mright\} = \mleft\{ c \mright\} \lor \mleft\{ a \mright\} = \mleft\{ c, d \mright\} \label{eq:for-curly-a}\\
        \mleft\{ a, b \mright\} = \mleft\{ c \mright\} \lor \mleft\{ a, b \mright\} = \mleft\{ c, d \mright\} \label{eq:for-curly-ab}
    \end{gather}
    Since \( a \neq b \), necessarily \( \mleft\{ a, b \mright\} \neq \mleft\{ c \mright\} \). Therefore, \( \mleft\{ a, b \mright\} = \mleft\{ c, d \mright\} \)
    must hold by \eqref{eq:for-curly-ab}. Now suppose that \( \mleft\{ a \mright\} = \mleft\{ c, d \mright\} \). But this would imply \( a = c = d \). But
    \( \mleft\{ a, b \mright\} = \mleft\{ c, d \mright\} \). This means that \( \mleft\{ c, b \mright\} = \mleft\{ c \mright\} \), so that \( a = c = b \),
    a contradiction. Thus, \( \mleft\{ a \mright\} \neq \mleft\{ c, d \mright\} \). By applying
    \aref{rule-of-inference:disjunctive-syllogism}{Disjunctive Syllogism} to \eqref{eq:for-curly-a}, we conclude that \( \mleft\{ a \mright\} = \mleft\{ c \mright\} \).
    \begin{equation*}
        \therefore \, a = c
    \end{equation*}
    With that, \( b = d \) follows immediately from the fact that \( \mleft\{ a, b \mright\} = \mleft\{ c, d\mright\} \).

    \vspace{0.5cm}
    \noindent
    In both cases, \( a = c \land b = d \) is implied. Thus,
    \begin{equation}
        \label{eq:ordered-pair-equality-forward-implication}
        \orderedpair{a}{b} = \orderedpair{c}{d} \implies \mleft( a = c \land b = d \mright)
    \end{equation}
    Proving the converse is trivial. Suppose that \( a = c \) and \( b = d \) for arbitrary ordered pairs \( \alpha = \orderedpair{a}{b} \)
    and \( \beta = \orderedpair{c}{d} \). Then \( \alpha = \beta \) is necessarily true because \( \mleft\{ a \mright\} = \mleft\{ c \mright\} \)
    and \( \mleft\{ a, b \mright\} = \mleft\{ c, d \mright\} \). Thus,
    \begin{equation}
        \label{eq:ordered-pair-equality-backward-implication}
        \mleft( a = c \land b = d \mright) \implies \orderedpair{a}{b} = \orderedpair{c}{d}
    \end{equation}
    Together, \eqref{eq:ordered-pair-equality-forward-implication} and \eqref{eq:ordered-pair-equality-backward-implication} establish
    the required biconditional.
\end{proof}

\begin{thm}
    \label{thm:ordered-pair-belongs-to-the-power-set-of-the-powerset}
    If \( a \in A \) and \( b \in A \), then \( \orderedpair{a}{b} \in \powerset{\powerset{A}} \).
    \begin{equation*}
        \forall a,b \, \forall A \, \mleft( a \in A \land b \in A \implies \orderedpair{a}{b} \in \powerset{\powerset{A}} \mright)
    \end{equation*}
\end{thm}

\begin{proof}
    Let \( A \) be an arbitrary set. Suppose that \( a \in A \) and \( b \in A \). Clearly, \( \mleft\{ a \mright\} \subseteq A \)
    and \( \mleft\{ a, b \mright\} \subseteq A \). Therefore, \( \mleft\{ a \mright\} \in \powerset{A} \) and
    \( \mleft\{ a, b \mright\} \in \powerset{A} \). Now let \( x_{0} \in \orderedpair{a}{b}\)
    be arbitrary. By \cref{dfn:kuratowski-ordered-pairs}, it follows that,
    \begin{equation*}
        x_{0} = \mleft\{ a \mright\} \lor x_{0} = \mleft\{ a, b \mright\}
    \end{equation*}
    Since \( \mleft\{ a \mright\} \in \powerset{A} \) and \( \mleft\{ a, b \mright\} \in \powerset{A} \) are both true, it follows
    by \aref{rule-of-inference:disjunction-elimination}{Disjunction Elimination} that \( x_{0} \in \powerset{A} \).
    In other words, \( x_{0} \in \orderedpair{a}{b} \implies x_{0} \in \powerset{A} \)
    for any arbitrary \( x_{0} \). Thus, by \aref{rule-of-inference:universal-generalization}{Universal Generalization},
    \begin{equation*}
        \forall x \, \mleft( x \in \orderedpair{a}{b} \implies x \in \powerset{A} \mright)
    \end{equation*}
    By \cref{dfn:subset,dfn:power-set}, this means that \( \orderedpair{a}{b} \subseteq \powerset{A} \) and \( \orderedpair{a}{b} \in \powerset{\powerset{A}} \).
    \begin{equation*}
        \therefore \,
        a \in A \land b \in A
        \implies
        \orderedpair{a}{b} \in \powerset{\powerset{A}}
    \end{equation*}
    Since \( a \), \( b \), \( A \) are all arbitrary, the theorem follows by \aref{rule-of-inference:universal-generalization}{Universal Generalization}.
\end{proof}

\begin{thm}
    \label{thm:the-coordinates-of-an-orderedpair-belongs-to-the-union-of-the-union-of-the-set-containing-the-orderpair}
    If \( \orderedpair{a}{b} \in A \), then \( a \in \Union{\Union{A}} \land b \in \Union{\Union{A}} \).
    \begin{equation*}
        \forall a,b \, \forall A \, \mleft( \orderedpair{a}{b} \in A \implies
        a \in \Union{\Union{A}} \land b \in \Union{\Union{A}}\mright)
    \end{equation*}
\end{thm}

\begin{proof}
    Let \( A \) be an arbitrary set and \orderedpair{a}{b} be an arbitrary ordered pair.
    Suppose that \( \orderedpair{a}{b} \in A \).
    Since \( \mleft\{ a, b \mright\} \in \orderedpair{a}{b} \) by \cref{dfn:kuratowski-ordered-pairs}
    and \( \orderedpair{a}{b} \in A \), it follows that,
    \begin{equation*}
        \exists S  \mleft( S \in A \land \mleft\{ a, b \mright\} \in S \mright)
    \end{equation*}
    Then, by \cref{dfn:union-of-sets}, it follows that \( \mleft\{ a, b \mright\} \in \Union{A} \).
    But \( a \in \mleft\{ a, b \mright\} \) and \( b \in \mleft\{ a, b \mright\} \), so we have,
    \begin{gather*}
        \exists S  \mleft( S \in \Union{A} \land a \in S \mright) \\
        \exists S  \mleft( S \in \Union{A} \land b \in S \mright)
    \end{gather*}
    Again, by \cref{dfn:union-of-sets}, we have \( a \in \Union{\Union{A}} \) and
    \( b \in \Union{\Union{A}} \).
\end{proof}

\newpage
\section{Sets of Ordered Pairs}
Having rigorously defined ordered pairs as sets, the obvious next step would be to
collect them together into sets---that is, to form sets of ordered pairs.
However, whenever one attempts to introduce new sets, it is crucial to first establish
their existence on the basis of the axioms of set theory. This is the subject of the
next two theorems, which describe the sufficient conditions for the existence of sets of ordered pairs.

\begin{thm}
    \label{thm:sufficient-condition-for-the-existence-of-sets-of-ordered-pairs-1}
    Given any binary predicate \( \varphi \),
    \begin{equation*}
        \exists A \, \forall x,y \, \Big( \varphi(x,y) {\implies} \orderedpair{x}{y} \in A \Big)
        {\implies}
        \exists! S \, \forall x \, \Big( x \in S {\iff} \exists a,b \, \big[ x = \orderedpair{a}{b} \land \varphi(a,b) \big] \Big)
    \end{equation*}
\end{thm}

\begin{proof}
    Let \( \varphi \) be a binary predicate\,\footnote{A predicate with two free variables.} such that,
    \begin{equation}
        \label{eq:phi-guarantees-ordered-pair-belonging-to-an-existing-set}
        \forall x,y \, \Big( \varphi(x,y) \implies \orderedpair{x}{y} \in A \Big)
    \end{equation}
    for some set \( A \). Now define a unary predicate \( \Gamma \) as follows,
    \begin{equation}
        \label{eq:define-gamma}
        \Gamma(x) \coloneq
        \exists a,b \, \Big( x = \orderedpair{a}{b} \land \varphi(a, b) \Big)
    \end{equation}
    Let \( x_{0} \) be arbitrary object such that \( \Gamma(x_{0}) \) holds. Thus, by
    \eqref{eq:define-gamma}, \( x_{0} = \orderedpair{a}{b} \land \varphi(a, b) \)
    for some \( a \) and \( b \). Since \( \varphi(a,b) \), \( \orderedpair{a}{b} \in A \)
    by \eqref{eq:phi-guarantees-ordered-pair-belonging-to-an-existing-set}.
    Since \( x_{0} = \orderedpair{a}{b} \), it follows that \( x_{0} \in A \).
    \begin{equation*}
        \therefore \, \Gamma(x_{0}) \implies x_{0} \in A
    \end{equation*}
    Since \( x_{0} \) is arbitrary, we have,
    \begin{equation*}
        \forall x \, \big( \Gamma(x) \implies x \in A \big)
    \end{equation*}
    By \cref{thm:sufficient-condition-for-unrestricted-comprehension}, there exists a unique
    set \( S \) such that,
    \begin{equation*}
        S = \setbuilderbig{x}{\Gamma(x)}
          = \setbuilderbig{x}{\exists a,b \, \big[ x = \orderedpair{a}{b} \land \varphi(a,b) \big]} \qedhere
    \end{equation*}
\end{proof}

\begin{thm}
    \label{thm:sufficient-condition-for-the-existence-of-sets-of-ordered-pairs-2}
    Given any binary predicate \( \varphi \),
    \begin{multline*}
        \exists A,B \, \forall x,y \, \Big( \varphi(x,y) \implies x \in A \land y \in B \Big)
        \implies \\
        \exists! S \, \forall x \, \Big( x \in S \iff \exists a,b \, \big[ x = \orderedpair{a}{b} \land \varphi(a,b) \big] \Big)
    \end{multline*}
\end{thm}

\begin{proof}
    Let \( \varphi \) be any binary predicate. Suppose that,
    \begin{equation}
        \label{eq:phi-xy-only-if-x-in-A-and-y-in-B}
        \forall x,y \, \Big( \varphi(x,y) \implies x \in A \land y \in B \Big)
    \end{equation}
    for some \( A \) and \( B \). Let \( x_{0} \) and \( y_{0} \) be arbitrary objects
    such that \( \varphi(x_{0}, y_{0}) \) is satisfied. By \eqref{eq:phi-xy-only-if-x-in-A-and-y-in-B},
    it follows that \( x_{0} \in A \) and \( y_{0} \in B \). Thus, \( x_{0} \in A \union B \) and
    \( y_{0} \in A \union B\),
    so necessarily \( \orderedpair{x_{0}}{y_{0}} \in \powerset{\powerset{A \union B}} \)
    by \cref{thm:ordered-pair-belongs-to-the-power-set-of-the-powerset}.
    \begin{equation*}
        \therefore \, \varphi(x_{0}, y_{0}) \implies \orderedpair{x_{0}}{y_{0}} \in \powerset{\powerset{A \union B}}
    \end{equation*}
    Since \( x_{0} \) and \( y_{0} \) are arbitrary, we conclude by \aref{rule-of-inference:universal-generalization}{Universal Generalization} that,
    \begin{equation*}
        \forall x,y \, \Big( \varphi(x,y) \implies \orderedpair{x}{y} \in \powerset{\powerset{A \union B}} \Big)
    \end{equation*}
    The theorem follows immediately from \cref{thm:sufficient-condition-for-the-existence-of-sets-of-ordered-pairs-1}.
\end{proof}

\begin{dfn}[Set-Builder Notation]
    \label{dfn:set-builder-notation-ordered-pair}
    For any binary predicate \( \varphi \), \setbuilder{\orderedpair{x}{y}}{\varphi(x)} denotes the set
    of all ordered pairs whose coordinates satisfy \( \varphi \). In other words,
    \begin{equation*}
        \forall x \, \Big( x \in \big\{ \orderedpair{a}{b} \mid \varphi(a,b) \big\} \iff
        \exists a,b \, \big[ x = \orderedpair{a}{b} \land \varphi(a,b) \big] \Big)
    \end{equation*}
\end{dfn}

\begin{thm}
    \label{thm:existence-of-the-domain-of-a-set-of-ordered-pairs}
    If \( R \) is a set of ordered pairs, then the set
    \setbuilderbig{x}{\exists y \, \mleft( \orderedpair{x}{y} \in R \mright)} exists.
\end{thm}

\begin{proof}
    Let \( R \) be an arbitrary set of ordered pairs. Let the predicate \( \varphi \) be defined as,
    \begin{equation*}
        \varphi(x) \coloneq \exists y \, \big( \orderedpair{x}{y} \in R \big)
    \end{equation*}
    Let \( x_{0} \) be an arbitrary object such that \( \varphi(x_{0}) \) holds. Thus, by the definition of \( \varphi \),
    \( \orderedpair{x_{0}}{y_{0}} \in R \) for some \( y_{0} \).
    By \cref{thm:the-coordinates-of-an-orderedpair-belongs-to-the-union-of-the-union-of-the-set-containing-the-orderpair},
    it follows that \( x_{0} \in \Union{\Union{R}} \). Thus,
    \begin{equation*}
        \varphi(x_{0}) \implies x_{0} \in \Union{\Union{R}}
    \end{equation*}
    Since \( x_{0} \) is arbitrary,
    \begin{equation*}
        \forall x \, \mleft( \varphi(x) \implies x \in \Union{\Union{R}} \mright)
    \end{equation*}
    By \cref{thm:sufficient-condition-for-unrestricted-comprehension},
    the set \( \setbuilderbig{x}{\varphi(x)} = \setbuilderbig{x}{\exists y \, \mleft( \orderedpair{x}{y} \in R \mright)} \)
    exists.
\end{proof}

\begin{dfn}[Domain]
    \label{dfn:domain-of-a-set-of-ordered-pairs}
    If \( R \) is a set of ordered pairs, then the \keyword{domain} of \( R \) is the set given by,
    \begin{equation*}
        \dom R \coloneq \setbuilderbig{x}{\exists y \, \big(\orderedpair{x}{y} \in R\big)}
    \end{equation*}
\end{dfn}

\begin{thm}
    \label{thm:existence-of-the-range-of-a-set-of-ordered-pairs}
    If \( R \) is a set of ordered pairs, then the set
    \setbuilderbig{y}{\exists x \, \mleft(\orderedpair{x}{y} \in R\mright)} exists.
\end{thm}

\begin{proof}
    Let \( R \) be an arbitrary set of ordered pairs. Let the predicate \( \varphi \) be defined as,
    \begin{equation*}
        \varphi(y) \coloneq \exists x \, \big( \orderedpair{x}{y} \in R \big)
    \end{equation*}
    Let \( y_{0} \) be an arbitrary object such that \( \varphi(y_{0}) \) is satisfied. By the definition
    of \( \varphi \), this means that \( \orderedpair{x_{0}}{y_{0}} \in R \) for some \( x_{0} \). By
    \cref{thm:the-coordinates-of-an-orderedpair-belongs-to-the-union-of-the-union-of-the-set-containing-the-orderpair},
    it follows that \( y_{0} \in \Union{\Union{f}} \). Thus,
    \begin{equation*}
        \varphi(y_{0}) \implies y_{0} \in \Union{\Union{R}}
    \end{equation*}
    Since \( y_{0} \) is arbitrary, \aref{rule-of-inference:universal-generalization}{Universal Generalization} yields,
    \begin{equation*}
        \forall y \, \mleft( \varphi(y) \implies y \in \Union{\Union{R}} \mright)
    \end{equation*}
    By \cref{thm:sufficient-condition-for-unrestricted-comprehension}, the set
    \setbuilder{y}{\exists x \, \mleft( \orderedpair{x}{y} \in R \mright)} exists.
\end{proof}

\begin{dfn}[Range]
    \label{dfn:range-of-a-set-of-ordered-pairs}
    If \( R \) is a set of ordered pairs, then the \keyword{range} of \( R \) is the set given by,
    \begin{equation*}
        \ran R \coloneq \setbuilderbig{y}{\exists x \, \mleft( \orderedpair{x}{y} \in R \mright)}
    \end{equation*}
\end{dfn}

\begin{thm}
    \phantomsection \label{thm:existence-of-the-image-of-a-set-of-ordered-pairs-under-another-set}
    If \( R \) is a set of ordered pairs and \( A \) is an arbitrary set,
    then the set \( \setbuilder{y}{\exists x \, \mleft( x \in A \,\,{\land}\,\, \orderedpair{x}{y} \in R \mright)} \) exists.
\end{thm}

\begin{proof}
    Let the predicate \( \varphi \) be defined as follows:
    \begin{equation*}
        \varphi(y) \coloneq \exists x \, \big[ x \in A \,\,{\land}\,\, \orderedpair{x}{y} \in R \big]
    \end{equation*}
    Let \( y_{0} \) be an arbitrary object such that \( \varphi(y_{0}) \) holds. Thus, \( \orderedpair{x_{0}}{y_{0}} \in R \) for some \( x_{0} \in A \).
    But this implies that \( \exists x \, \mleft[ \orderedpair{x}{y_{0}} \in R \mright] \).
    In other words, \( y_{0} \in \ran R \) by \cref{dfn:range-of-a-set-of-ordered-pairs}.
    By the \aref{rule-of-inference:deduction-theorem}{Deduction Theorem}, we have \( \varphi(y_{0}) {\implies} y_{0} \in \ran R \) for any arbitrary \( y_{0} \).
    \aref{rule-of-inference:universal-generalization}{Universal Generalization} then yields,
    \begin{equation*}
        \forall y \, \big[ \varphi(y) \implies y \in \ran R \big]
    \end{equation*}
    Since \( \ran R \) exists by \cref{thm:existence-of-the-range-of-a-set-of-ordered-pairs}, we thus have,
    \begin{equation*}
        \exists A \, \forall y \, \big[ \varphi(y) \implies y \in A \big]
    \end{equation*}
    The existence of \( \setbuilder{y}{\varphi(y)} \) follows immediately from \cref{thm:sufficient-condition-for-unrestricted-comprehension}.
\end{proof}

\begin{dfn}[Image]
    \label{dfn:image-of-a-set-of-ordered-pairs}
    For any set \( A \) and any set of ordered pairs \( R \), the \keyword{image} of \( A \) under \( R \)
    is the set given by,
    \begin{equation*}
        \imageSet{R}{A} \coloneq \setbuilder{y \in \ran R}{\exists x \, \big( x \in A \,\,{\land}\,\, \orderedpair{x}{y} \in R \big)}
    \end{equation*}
\end{dfn}

\begin{thm}
    \label{thm:existence-of-the-inverse-of-a-set-of-ordered-pairs}
    For any set of ordered pairs \( R \), the set \setbuilderbig{\orderedpair{y}{x}}{\orderedpair{x}{y} \in R} exists.
\end{thm}

\begin{proof}
    Let \( R \) be an arbitrary set of ordered pairs.
    By \cref{thm:existence-of-the-domain-of-a-set-of-ordered-pairs,thm:existence-of-the-range-of-a-set-of-ordered-pairs}, \( \dom R \)
    and \( \ran R \) exist and are given by,
    \begin{align}
        \dom R &=
        \setbuilderbig{x}{\exists y \, \big( \orderedpair{x}{y} \in R \big)} \label{eq:domain-of-f-set-builder-form}\\
        \ran R &=
        \setbuilderbig{y}{\exists x \, \big( \orderedpair{x}{y} \in R \big)} \label{eq:range-of-f-set-builder-form}
    \end{align}
    Let the binary predicate \( \varphi \) be defined as,
    \begin{equation*}
        \varphi(y, x) \coloneq \orderedpair{x}{y} \in R
    \end{equation*}
    Let \( y_{0} \) and \( x_{0} \) be arbitrary objects such that \( \varphi(y_{0}, x_{0}) \) is satisfied. Thus, by
    the definition of \( \varphi \), it follows that \( \orderedpair{x_{0}}{y_{0}} \in R \), which implies,
    \begin{gather*}
        \exists x \, \big( \orderedpair{x}{y_{0}} \in R \big) \\
        \exists y \, \big( \orderedpair{x_{0}}{y} \in R \big)
    \end{gather*}
    Thus, \( y_{0} \in \ran R \) and \( x_{0} \in \dom f \) by \eqref{eq:range-of-f-set-builder-form}
    and \eqref{eq:domain-of-f-set-builder-form}. Clearly,
    \begin{equation*}
        y_{0} \in \mleft( \dom R \union \ran R \mright)
        \land
        x_{0} \in \mleft( \dom R \union \ran R \mright)
    \end{equation*}
    Since \( y_{0} \) and \( x_{0} \) are arbitrary, \aref{rule-of-inference:universal-generalization}{Universal Generalization} yields,
    \begin{equation*}
        \forall y \, \forall x \, \big( \varphi(y, x) \implies y \in \mleft( \dom R \union \ran R \mright)
        \land x \in \mleft( \dom R \union \ran R \mright) \big)
    \end{equation*}
    Thus, the set \( \setbuilderbig{\orderedpair{y}{x}}{\varphi(y,x)} = \setbuilderbig{\orderedpair{y}{x}}{\orderedpair{x}{y} \in R} \)
    exists by \cref{thm:sufficient-condition-for-the-existence-of-sets-of-ordered-pairs-2}.
\end{proof}

\begin{dfn}[Inverse]
    \label{dfn:inverse-of-a-set-of-ordered-pairs}
    For any set of ordered pairs \( R \), the \keyword{inverse} of \( R \) is the set given by,
    \begin{equation*}
        R^{-1} \coloneq \setbuilderbig{\orderedpair{y}{x}}{\orderedpair{x}{y} \in R}
    \end{equation*}
\end{dfn}

\begin{thm}
    \label{thm:existence-of-the-composition-of-two-sets-of-ordered-pairs}
    If \( R \) and \( S \) are sets of ordered pairs, then the set,
    \begin{equation*}
        \setbuilderbig{\orderedpair{x}{y}}{\exists z \, \big( \orderedpair{x}{z} \in S \land \orderedpair{z}{y} \in R \big)}
    \end{equation*}
    exists.
\end{thm}

\begin{proof}
    Let \( R \) and \( S \) be arbitrary sets of ordered pairs. Let the binary predicate \( \varphi \) be defined as,
    \begin{equation*}
        \varphi(x,y) \coloneq \exists z \, \big( \orderedpair{x}{z} \in S \land \orderedpair{z}{y} \in R \big)
    \end{equation*}
    Let \( x_{0} \) and \( y_{0} \) be arbitrary objects such that \( \varphi(x_{0}, y_{0}) \) is satisfied. Thus,
    \( \orderedpair{x_{0}}{z_{0}} \in S \) and \( \orderedpair{z_{0}}{y_{0}} \in R \) for some \( z_{0} \).
    \begin{equation*}
        \therefore \, \exists z \, \big( \orderedpair{x_{0}}{z} \in S \big) \land \exists z \, \big( \orderedpair{z}{y_{0}} \in R \big)
    \end{equation*}
    It follows from \cref{dfn:domain-of-a-set-of-ordered-pairs,dfn:range-of-a-set-of-ordered-pairs} that
    \( x_{0} \in \dom S \) and \( y_{0} \in \ran R \). In other words, for any \( x_{0} \) and \( y_{0} \),
    \begin{equation*}
        \varphi(x_{0}, y_{0}) \implies x_{0} \in \dom S \land y_{0} \in \ran R
    \end{equation*}
    By \aref{rule-of-inference:universal-generalization}{Universal Generalization}, we have,
    \begin{equation*}
        \forall x,y \, \Big( \varphi(x,y) \implies x \in \dom S \land y \in \ran R \Big)
    \end{equation*}
    Thus, the set \(\setbuilderbig{\orderedpair{x}{y}}{\exists z \, \big( \orderedpair{x}{y} \in S \land \orderedpair{y}{z} \in R \big)} \)
    exists by \cref{thm:sufficient-condition-for-the-existence-of-sets-of-ordered-pairs-2}.
\end{proof}

\begin{dfn}[Composition]
    \label{dfn:composition-of-two-sets-of-ordered-pairs}
    If \( R \) and \( S \) are sets of ordered pairs, then the \keyword{composition} of \( S \) and \( R \) is the set,
    \begin{equation*}
        \composition{R}{S} \coloneq \setbuilderbig{\orderedpair{x}{y}}{\exists z \, \big( \orderedpair{x}{z} \in S \land \orderedpair{z}{y} \in R \big)}
    \end{equation*}
\end{dfn}

\subsection{Properties of Sets of Ordered Pairs} \label{ssec:properties-of-sets-of-ordered-pairs}
\begin{thm}
    \label{thm:empty-image-of-a-set-of-ordered-pairs}
    For any set of ordered pairs \( R \), \( \imageSet{R}{\emptyset} = \emptyset\).
\end{thm}

\begin{proof}
    Let \( R \) be an arbitrary set of ordered pairs. By \cref{dfn:image-of-a-set-of-ordered-pairs},
    \begin{equation*}
        \imageSet{R}{\emptyset} = \setbuilderBig{y \in \ran R}{\exists x \, \big( x \in \emptyset \land \orderedpair{x}{y} \in R \big)}
    \end{equation*}
    Let the predicate \( \varphi \) be defined as,
    \begin{equation}
        \label{eq:let-phi-represent-the-predicate-for-the-image-of-R-under-empty}
        \varphi(y) \coloneq \exists x \, \big( x \in \emptyset \land \orderedpair{x}{y} \in R \big)
    \end{equation}
    Thus,
    \begin{equation}
        \label{eq:restating-the-image-of-R-under-empty-in-terms-of-phi}
        \forall y \, \Big( y \in \imageSet{R}{\emptyset}
        \iff
        y \in \ran R
        \land
        \varphi(y)
        \Big)
    \end{equation}
    Using the laws of elementary logic, we can re-express formula \eqref{eq:let-phi-represent-the-predicate-for-the-image-of-R-under-empty}
    equivalently as,
    \begin{align}
        \label{eq:phi-reformulated-equivalently}
        \varphi(y) &\iff \exists x \, \big( x \in \emptyset \land \orderedpair{x}{y} \in R \big) \notag \\
        &\iff
        \lnot \, \lnot \, \exists x \, \big( x \in \emptyset \land \orderedpair{x}{y} \in R \big) \notag \\
        &\iff
        \lnot \, \forall x \, \big( x \notin \emptyset \lor \orderedpair{x}{y} \notin R \big) \notag \\
        &\iff
        \lnot \, \forall x \, \big( x \in \emptyset \implies \orderedpair{x}{y} \notin R \big)
    \end{align}
    Let \( y_{0} \in \ran R \) be arbitrary. Let \( x_{0} \) be an arbitrary object. By \cref{dfn:empty-set},
    clearly \( x_{0} \notin \emptyset \). Thus, the formula \( x_{0} \in \emptyset \implies \orderedpair{x_{0}}{y_{0}} \notin R \)
    is vacuously true. Since \( x_{0} \) is arbitrary, we infer by \aref{rule-of-inference:universal-generalization}{Universal Generalization},
    \begin{equation*}
        \forall x \, \big( x \in \emptyset \implies \orderedpair{x}{y_{0}} \notin R \big)
    \end{equation*}
    Thus, \( \lnot \, \varphi(y_{0}) \) by \eqref{eq:phi-reformulated-equivalently}. Furthermore, \( y_{0} \notin \imageSet{R}{\emptyset} \)
    by \eqref{eq:restating-the-image-of-R-under-empty-in-terms-of-phi}. Since \( y_{0} \) is arbitrary, we conclude by \aref{rule-of-inference:universal-generalization}{Universal Generalization}
    that,
    \begin{equation*}
        \forall y \, \big( y \notin \imageSet{R}{\emptyset} \big)
    \end{equation*}
    By \cref{dfn:empty-set}, it follows that \( \imageSet{R}{\emptyset} = \emptyset \).
\end{proof}

\begin{thm}
    \label{thm:inverse-of-the-inverse-of-a-set-of-ordered-pairs-is-the-set-itself}
    For any set of ordered pairs \( R \), \( R = \mleft( R^{-1} \mright)^{-1} \)
\end{thm}

\begin{proof}
    Let \( R \) be an arbitrary set of ordered pairs. \cref{dfn:inverse-of-a-set-of-ordered-pairs} yields the following equalities:
    \begin{equation*}
        \mleft( R^{-1} \mright)^{-1}
        = \setbuilderbig{\orderedpair{x}{y}}{\orderedpair{y}{x} \in R^{-1}}
        = \setbuilderbig{\orderedpair{x}{y}}{\orderedpair{x}{y} \in R}
        = R \qedhere
    \end{equation*}
\end{proof}

\begin{thm}
    \label{thm:domain-is-the-range-of-inverse-and-vice-versa}
    If \( R \) is a set of ordered pairs, then,
    \begin{gather*}
        \dom R = \ran R^{-1} \\
        \ran R = \dom R^{-1}
    \end{gather*}
\end{thm}

\begin{proof}
    Let \( R \) be an arbitrary set of ordered pairs. By \cref{dfn:inverse-of-a-set-of-ordered-pairs}, \( R^{-1} \)
    is given by,
    \begin{equation*}
        R^{-1} = \setbuilderbig{\orderedpair{y}{x}}{\orderedpair{x}{y} \in R}
    \end{equation*}
    Let \( x_{0} \) be an arbitrary object. By \cref{dfn:range-of-a-set-of-ordered-pairs}, we have,
    \begin{align*}
        x_{0} \in \ran R^{-1}
        &\iff
        \exists y \, \big( \orderedpair{y}{x_{0}} \in R^{-1} \big)\\
        &\iff
        \exists y \, \big( \orderedpair{x_{0}}{y} \in R \big)\\
        &\iff
        x_{0} \in \dom R
    \end{align*}
    Since \( x_{0} \) is arbitrary, we have,
    \begin{equation*}
        \forall x \, \mleft( x \in \dom R \iff x \in \ran R^{-1} \mright)
    \end{equation*}
    Thus, \( \dom R = \ran R^{-1} \) follows immediately from the \nameref{axm:axiom-of-extensionality}.

    \vspace{0.3cm}
    \noindent
    The second equality can be proven the same way. Let \( y_{0} \) be an arbitrary object
    By \cref{dfn:domain-of-a-set-of-ordered-pairs}, we have,
    \begin{align*}
        y_{0} \in \dom R^{-1}
        &\iff
        \exists x \, \big( \orderedpair{y_{0}}{x} \in R^{-1} \big)\\
        &\iff
        \exists x \, \big( \orderedpair{x}{y_{0}} \in R \big)\\
        &\iff
        y_{0} \in \ran R
    \end{align*}
    By \aref{rule-of-inference:universal-generalization}{Universal Generalization} and the \nameref{axm:axiom-of-extensionality}, we conclude,
    \begin{equation*}
        \dom R^{-1} = \ran R \qedhere
    \end{equation*}
\end{proof}

\begin{thm}
    \label{thm:image-under-a-union-is-equal-to-the-union-of-individual-images}
    If \( R \) is a set of ordered pairs and \( S \) is a non-empty system of sets, then,
    \begin{equation*}
        \imageSet{R}{\Union S} = \Union \setbuilderBig{t}{\exists z \, \mleft( z \in S \land t = \imageSet{R}{z} \mright)}
    \end{equation*}
\end{thm}

\begin{proof}
    Let \( R \) be a set of ordered pairs. Let \( S \) be a non-empty system of sets. Let the set \( T \) be defined as,
    \begin{equation}
        \label{eq:T-as-the-set-of-all-images-under-R}
        T = \setbuilderBig{t}{\exists z \, \big( z \in S \land t = \imageSet{R}{z} \big)}
    \end{equation}
    We first prove \( T \)'s existence. Let \( t_{0} \) be an arbitrary object such that,
    \begin{equation*}
        \exists z \, \big( z \in S \land t_{0} = \imageSet{R}{z} \big)
    \end{equation*}
    Therefore, we have \( t_{0} = \imageSet{R}{z_{0}} \) for some \( z_{0} \in S \). By \cref{dfn:image-of-a-set-of-ordered-pairs},
    it follows that \( \imageSet{R}{z_{0}} \subseteq \ran R \). Thus, \( t_{0} \subseteq \ran R \), which in turn implies
    \( t_{0} \in \powerset{\ran R} \).
    \begin{equation*}
        \therefore \, \forall t \, \Big( \exists z \, \big( z \in S \land t = \imageSet{R}{z} \big) \implies t \in \powerset{\ran R} \Big)
    \end{equation*}
    Thus, \( T \) exists by \cref{thm:sufficient-condition-for-unrestricted-comprehension}.

    \vspace{0.3cm}
    \noindent
    Naturally, the main goal is to prove that,
    \begin{equation*}
        \imageSet{R}{\Union S} = \Union T
    \end{equation*}
    To do so, we first let \( y_{0} \) represent an arbitrary object and then apply
    \aref{rule-of-inference:universal-instantiation}{Universal Instantiation} on \cref{dfn:image-of-a-set-of-ordered-pairs} to obtain,
    \begin{equation}
        \label{eq:initial-statement-for-the-image-of-union-S-under-R}
        y_{0} \in \imageSet{R}{\Union S}
        \iff
        \exists x \, \mleft( x \in \Union S \land \orderedpair{x}{y_{0}} \in R \mright)
    \end{equation}
    The laws of elementary logic (\cref{tautology:grouping-conjuncts-into-existentially-quantified-formula})
    can then be applied to transform \eqref{eq:initial-statement-for-the-image-of-union-S-under-R}
    into a series of equivalent formulas:
    \begin{align}
        \label{eq:final-statement-for-the-image-of-union-S-under-R}
        y_{0} \in \imageSet{R}{\Union S}
        &\iff
        \exists x \, \mleft( x \in \Union S \land \orderedpair{x}{y_{0}} \in R \mright) \notag \\
        &\iff
        \exists x \, \Big( \exists z \, \mleft( z \in S \land x \in z \mright) \land \orderedpair{x}{y_{0}} \in R \Big) \notag \\
        &\iff
        \exists x \, \exists z \, \Big( z \in S \land x \in z \land \orderedpair{x}{y_{0}} \in R \Big) \notag \\
        &\iff
        \exists z \, \exists x \, \Big( z \in S \land x \in z \land \orderedpair{x}{y_{0}} \in R \Big) \notag \\
        &\iff
        \exists z \, \Big( z \in S \land \exists x \, \big( x \in z \land \orderedpair{x}{y_{0}} \in R \big) \Big) \notag \\
        &\iff
        \exists z \, \Big( z \in S \land y_{0} \in \imageSet{R}{z} \Big)
    \end{align}
    We then repeat the process for the right-hand side of the equality. Start by applying \cref{dfn:union-of-sets}
    and \eqref{eq:T-as-the-set-of-all-images-under-R} to \( \Union T \).
    \begin{equation*}
        y_{0} \in \Union T
        \iff
        \exists t \, \mleft( t \in T \land y_{0} \in t \mright)
        \iff
        \exists t \, \Big( \exists z \, \big( z \in S \land t = \imageSet{R}{z} \big) \land y_{0} \in t \Big)
    \end{equation*}
    Using \cref{tautology:grouping-conjuncts-into-existentially-quantified-formula}, the formula becomes,
    \begin{equation*}
        y_{0} \in \Union T
        \iff
        \exists z \, \Big( z \in S \land \exists t \, \mleft( t = \imageSet{R}{z} \land y_{0} \in t \mright) \Big)
    \end{equation*}
    By \cref{tautology:using-an-existentially-quantified-formula-to-express-the-fact-that-an-object-satisfies-a-formula}, we have,
    \begin{equation*}
        \exists t \, \big( t = \imageSet{R}{z} \land y_{0} \in t \big)
        \iff
        y_{0} \in \imageSet{R}{z}
    \end{equation*}
    Thus,
    \begin{equation}
        \label{eq:final-statement-of-the-union-of-T}
        y_{0} \in \Union T
        \iff
        \exists z \, \Big( z \in S \land y_{0} \in \imageSet{R}{z} \Big)
    \end{equation}
    Comparing \eqref{eq:final-statement-for-the-image-of-union-S-under-R} and \eqref{eq:final-statement-of-the-union-of-T},
    it is clear that,
    \begin{equation*}
        y_{0} \in \imageSet{R}{\Union S} \iff y_{0} \in \Union T
    \end{equation*}
    Since \( y_{0} \) is arbitrary, we conclude,
    \begin{equation*}
        \forall y \, \mleft( y \in \imageSet{R}{\Union S} \iff y \in \Union T \mright)
    \end{equation*}
    The theorem follows readily from the \nameref{axm:axiom-of-extensionality}.
\end{proof}

\begin{cor}
    \label{thm:image-of-A-union-B-under-R-is-the-union-of-the-image-of-A-under-R-and-the-image-of-B-under-R}
    Let \( R \) be a set of ordered pairs. Let \( A \) and \( B \) be two sets.
    Then \( \imageSet{R}{A \union B} = \imageSet{R}{A} \union \imageSet{R}{B} \).
\end{cor}

\begin{thm}
    \label{thm:image-of-an-intersection-is-a-subset-of-the-intersection-of-all-individual-images}
    If \( R \) is a set of ordered pairs and \( S \) is a non-empty system of sets, then,
    \begin{equation*}
        \imageSet{R}{\Intersect S} \subseteq \Intersect \setbuilderBig{t}{\exists z \, \big( z \in S \land t = \imageSet{R}{z} \big)}
    \end{equation*}
\end{thm}

\begin{proof}
    Let \( R \) be an arbitrary set of ordered pairs. Let \( S \) be a non-empty system of sets.
    Let \( T \) be the same set defined in \eqref{eq:T-as-the-set-of-all-images-under-R}. The goal
    is to prove that,
    \begin{equation*}
        \forall y \, \mleft( y \in \imageSet{R}{\Intersect S} \implies y \in \Intersect T \mright)
    \end{equation*}
    Let \( y_{0} \) represent an arbitrary object. By \eqref{eq:T-as-the-set-of-all-images-under-R} and
    \cref{dfn:intersection-of-sets}, it follows that,
    \begin{align}
        \label{eq:condition-for-y0-to-belong-to-the-intersection-of-T}
        y_{0} \in \Intersect T
        &\iff
        \forall t \, \Big( t \in T \implies y_{0} \in t \Big) \notag \\
        &\iff
        \forall t \, \Big( \exists z \, \big( z \in S \land t = \imageSet{R}{z} \big) \implies y_{0} \in t \Big) \notag \\
        &\iff
        \forall t \, \Big( \forall z \, \big( z \notin S \lor t \neq \imageSet{R}{z} \big) \lor y_{0} \in t \Big) \notag \\
        &\iff
        \forall t \, \forall z \, \Big( z \notin S \lor t \neq \imageSet{R}{z} \lor y_{0} \in t \Big) \notag \\
        &\iff
        \forall z \, \forall t \, \Big( z \notin S \lor t \neq \imageSet{R}{z} \lor y_{0} \in t \Big) \notag \\
        &\iff
        \forall z \, \Big( z \in S \implies \forall t \, \mleft( t = \imageSet{R}{z} \implies y_{0} \in t \mright) \Big) \notag \\
        &\iff
        \forall z \, \Big( z \in S \implies y_{0} \in \imageSet{R}{z} \Big)
    \end{align}
    Now suppose that \( y_{0} \in \imageSet{R}{\Intersect S} \).
    By \cref{dfn:image-of-a-set-of-ordered-pairs}, it follows that \( \orderedpair{x_{0}}{y_{0}} \in R \)
    for some \( x_{0} \in \Intersect S \).
    By \cref{dfn:intersection-of-sets}, necessarily,
    \begin{equation}
        \label{eq:x0-belongs-to-all-members-of-S}
        \forall z \, \mleft( z \in S \implies x_{0} \in z \mright)
    \end{equation}
    Let \( z_{0} \in S \) be arbitrary. By \eqref{eq:x0-belongs-to-all-members-of-S}, it follows that \( x_{0} \in z_{0} \). Thus, we have,
    \begin{gather*}
        z_{0} \in S \implies x_{0} \in z_{0} \land \orderedpair{x_{0}}{y_{0}} \in R \\
        \therefore \, z_{0} \in S \implies \exists x \, \big( x \in z_{0} \land \orderedpair{x}{y_{0}} \in R \big)
    \end{gather*}
    In other words, \( y_{0} \in \imageSet{R}{z_{0}} \) for any \( z_{0} \in S \). Since \( z_{0} \) is arbitrary, we have,
    \begin{equation*}
        \forall z \, \big( z \in S \implies y_{0} \in \imageSet{R}{z} \big)
    \end{equation*}
    Thus, \( y_{0} \in \Intersect T \) by \eqref{eq:condition-for-y0-to-belong-to-the-intersection-of-T}.
    \begin{equation*}
        \therefore \, y_{0} \in \imageSet{R}{\Intersect S} \implies y_{0} \in \Intersect T
    \end{equation*}
    Given that \( y_{0} \) is arbitrary, the theorem follows by \aref{rule-of-inference:universal-generalization}{Universal Generalization}.
\end{proof}

\begin{cor}
    \label{thm:the-image-of-A-intersect-B-under-R-is-a-subset-of-the-image-of-A-under-R-intersect-the-image-of-B-under-R}
    Let \( R \) be a set of ordered pairs. Let \( A \) and \( B \) be two sets.
    Then \( \imageSet{R}{A \intersect B} \subseteq \imageSet{R}{A} \intersect \imageSet{R}{B} \).
\end{cor}

\begin{thm}
    \label{thm:the-difference-of-image-is-a-subset-of-the-image-of-difference}
    Let \( R \) be a set of ordered pairs. Let \( A \) and \( B \) be two sets. Then \( \imageSet{R}{A} - \imageSet{R}{B} \subseteq \imageSet{R}{A - B} \)
\end{thm}

\begin{proof}
    Let \( y_{0} \) be an arbitrary object. We start by applying the definitions for \( \imageSet{R}{A} - \imageSet{R}{B} \) and \( \imageSet{R}{A-B} \).
    \begin{align}
        &y_{0} \in \imageSet{R}{A} - \imageSet{R}{B} \notag \\
        &\iff
        \exists x \, \big[ x \in A \land \orderedpair{x}{y_{0}} \in R \big]\land \lnot \, \exists x \, \big[ x \in B \land \orderedpair{x}{y_{0}} \in R \big] \notag \\
        &\iff
        \exists x \, \big[ x \in A \land \orderedpair{x}{y_{0}} \in R \big]
        \land
        \forall x \, \big[ x \in B \implies \orderedpair{x}{y_{0}} \notin R \big] \label{eq:y0-belongs-to-RA-minus-RB} \\
        \text{   } \notag \\
        &y_{0} \in \imageSet{R}{A-B} \notag \\
        &\iff
        \exists x \, \big[ x \in A \land x \notin B \land \orderedpair{x}{y_{0}} \in R \big] \label{eq:y0-belongs-to-R-A-minus-B}
    \end{align}
    Suppose that \( y_{0} \in \imageSet{R}{A} - \imageSet{R}{B} \). By \eqref{eq:y0-belongs-to-RA-minus-RB}, it follows that:
    \begin{equation*}
        \begin{aligned}
            \orderedpair{x_{0}}{y_{0}} \in R&
        \end{aligned}
        \text{  }
        \begin{aligned}
            &\text{for some \( x_{0} \in A \)}
        \end{aligned}
    \end{equation*}
    \aref{rule-of-inference:universal-instantiation}{Universal Instantiation} on the second conjunct of \eqref{eq:y0-belongs-to-RA-minus-RB} yields,
    \begin{equation*}
        x_{0} \in B \implies \orderedpair{x_{0}}{y_{0}} \notin R
    \end{equation*}
    Since \( \orderedpair{x_{0}}{y_{0}} \in R\), we conclude by \aref{rule-of-inference:modus-tollens}{Modus Tollens} that \( x_{0} \notin B \).
    \begin{equation*}
        \therefore \, \exists x \, \big[ x \in A \land x \notin B \land \orderedpair{x}{y_{0}} \in R \big]
    \end{equation*}
    By \eqref{eq:y0-belongs-to-R-A-minus-B}, it follows that \( y_{0} \in \imageSet{R}{A-B} \). By \aref{rule-of-inference:universal-generalization}{Universal Generalization}
    and the \aref{rule-of-inference:deduction-theorem}{Deduction Theorem},
    \begin{equation*}
        \forall y \, \Big[ y \in \imageSet{R}{A} - \imageSet{R}{B} \implies y \in \imageSet{R}{A-B} \Big]
    \end{equation*}
    Thus, \( \imageSet{R}{A} - \imageSet{R}{B} \subseteq \imageSet{R}{A-B} \).
\end{proof}

\begin{thm}
    \label{thm:the-symdiff-of-image-is-a-subset-of-the-image-of-symdiff}
    Let \( R \) be a set of ordered pairs. Let \( A \) and \( B \) be two sets.
    Then \( \imageSet{R}{A} \symdiff \imageSet{R}{B} \subseteq \imageSet{R}{A \symdiff B} \)
\end{thm}

\begin{proof}
    By \cref{set-identities:symmetric-difference-as-the-union-of-differences,thm:image-of-A-union-B-under-R-is-the-union-of-the-image-of-A-under-R-and-the-image-of-B-under-R},
    we have,
    \begin{align*}
        \imageSet{R}{A \symdiff B} &= \imageSet{R}{\mleft( A - B \mright) \union \mleft( B - A \mright)} = \imageSet{R}{A - B} \union \imageSet{R}{B - A} \\
        \imageSet{R}{A} \symdiff \imageSet{R}{B} &= \big(\imageSet{R}{A} - \imageSet{R}{B} \big) \union \big( \imageSet{R}{B} - \imageSet{R}{A} \big)
    \end{align*}
    By \cref{thm:the-difference-of-image-is-a-subset-of-the-image-of-difference}, \( \imageSet{R}{A} - \imageSet{R}{B} \subseteq \imageSet{R}{A - B} \)
    and \( \imageSet{R}{B} - \imageSet{R}{A} \subseteq  \imageSet{R}{B - A} \). And finally, by \cref{thm:the-union-of-subsets-is-a-subset-of-the-union},
    \begin{gather*}
        \big( \imageSet{R}{A} - \imageSet{R}{B} \big) \union \big( \imageSet{R}{B} - \imageSet{R}{A} \big)
        \subseteq
        \imageSet{R}{A - B} \union \imageSet{R}{B - A} \\
        \therefore \, \imageSet{R}{A} \symdiff \imageSet{R}{B} \subseteq \imageSet{R}{A \symdiff B} \qedhere
    \end{gather*}
\end{proof}

\begin{thm}
    \label{thm:A-subset-B-implies-image-of-A-under-R-subset-image-of-B-under-R}
    Let \( R \) be a set of ordered pairs. Let \( A \) and \( B \) be sets. Then,
    \begin{equation*}
        A \subseteq B \implies \imageSet{R}{A} \subseteq \imageSet{R}{B}
    \end{equation*}
\end{thm}

\begin{proof}
    Let \( y_{0} \) be an arbitrary object such that \( y_{0} \in \imageSet{R}{A} \). Thus
    \( \orderedpair{x_{0}}{y_{0}} \in R \) for some \( x_{0} \in A \) by \cref{dfn:image-of-a-set-of-ordered-pairs}.
    Since \( A \subseteq B \), \( x_{0} \in A \) implies \( x_{0} \in B \).
    \begin{equation*}
        \therefore \, \exists x \, \big( x \in B \land \orderedpair{x}{y_{0}} \in R \big)
    \end{equation*}
    By \cref{dfn:image-of-a-set-of-ordered-pairs}, \( y_{0} \in \imageSet{R}{B} \). In other words,
    \( y_{0} \in \imageSet{R}{A} {\implies} y_{0} \in \imageSet{R}{B} \) for any arbitrary \( y_{0} \).
    Thus, we have by \aref{rule-of-inference:universal-generalization}{Universal Generalization},
    \begin{equation*}
        \forall x \, \big( x \in \imageSet{R}{A} \implies x \in \imageSet{R}{B} \big)
    \end{equation*}
    By \cref{dfn:subset}, it follows that \( \imageSet{R}{A} \subseteq \imageSet{R}{B} \).
    And finally, we conclude by the \aref{rule-of-inference:deduction-theorem}{Deduction Theorem} that
    \( A \subseteq B \implies \imageSet{R}{A} \subseteq \imageSet{R}{B} \).
\end{proof}

\begin{thm}[Associativity]
    \label{thm:associativity-of-composition}
    Let \( R \), \( S \), and \( T \) be sets of ordered pairs. Then,
    \begin{equation*}
        \composition{\mleft( \composition{R}{S} \mright)}{T} = \composition{R}{\mleft( \composition{S}{T} \mright)}
    \end{equation*}
\end{thm}

\begin{proof}
    Let \( R \), \( S \), and \( T \) be sets of ordered pairs.
    Let \( x_{0} \) be an arbitrary object. By \cref{dfn:composition-of-two-sets-of-ordered-pairs},
    \begin{align}
        \label{eq:x0-belongs-to-lhs}
        &x_{0} \in \composition{\mleft( \composition{R}{S} \mright)}{T} \notag \\
        &\iff
        \exists a \, \exists b \,
        \Big[ x_{0} = \orderedpair{a}{b} \land \exists c \, \big( \orderedpair{a}{c} \in T \land \orderedpair{c}{b} \in \composition{R}{S} \big) \Big] \notag \\
        &\iff
        \exists a \, \exists b \,
        \Big[ x_{0} = \orderedpair{a}{b} \land \exists c \, \big( \orderedpair{a}{c} \in T \land
        \exists d \, \mleft( \orderedpair{c}{d} \in S \land \orderedpair{d}{b} \in R \mright) \big) \Big] \notag \\
        &\iff
        \exists a \, \exists b \,
        \Big[ x_{0} = \orderedpair{a}{b} \land \exists c \, \exists d \,
        \big( \orderedpair{a}{c} \in T \land \orderedpair{c}{d} \in S \land \orderedpair{d}{b} \in R \big) \Big] \notag \\
        &\iff
        \exists a \, \exists b \, \exists c \, \exists d \,
        \Big[ x_{0} = \orderedpair{a}{b} \land \orderedpair{a}{c} \in T \land \orderedpair{c}{d} \in S \land \orderedpair{d}{b} \in  R \Big]
    \end{align}
    Similarly,
    \begin{align}
        \label{eq:x0-belongs-to-rhs}
        &x_{0} \in \composition{R}{\mleft( \composition{S}{T} \mright)} \notag \\
        &\iff
        \exists a \, \exists b \,
        \Big[ x_{0} = \orderedpair{a}{b} \land \exists d \, \big( \orderedpair{a}{d} \in \composition{S}{T} \land \orderedpair{d}{b} \in R \big) \Big] \notag \\
        &\iff
        \exists a \, \exists b \,
        \Big[ x_{0} = \orderedpair{a}{b} \land \exists d \, \big( \exists c \,
        \mleft( \orderedpair{a}{c} \in T \land \orderedpair{c}{d} \in S \mright) \land \orderedpair{d}{b} \in R \big) \Big] \notag \\
        &\iff
        \exists a \, \exists b \, \exists c \,  \exists d \,
        \Big[ x_{0} = \orderedpair{a}{b} \land \orderedpair{a}{c} \in T \land \orderedpair{c}{d} \in S \land \orderedpair{d}{b} \in R \Big]
    \end{align}
    Comparing \eqref{eq:x0-belongs-to-lhs} and \eqref{eq:x0-belongs-to-rhs}, it is clear that,
    \begin{equation*}
        x_{0} \in \composition{ \mleft( \composition{R}{S} \mright) }{T} \iff x_{0} \in \composition{R}{\mleft( \composition{S}{T} \mright)}
    \end{equation*}
    Since \( x_{0} \) is arbitrary, we have by \aref{rule-of-inference:universal-generalization}{Universal Generalization},
    \begin{equation*}
        \forall x \, \Big( x \in \composition{\mleft( \composition{R}{S} \mright)}{T} \iff x \in \composition{R}{\mleft( \composition{S}{T} \mright)} \Big)
    \end{equation*}
    The theorem follows readily from the \nameref{axm:axiom-of-extensionality}.
\end{proof}

\begin{thm}
    \label{thm:inverse-of-composition}
    If \( R \) and \( S \) are sets of ordered pairs, then \( \mleft( \composition{R}{S} \mright)^{-1} = \composition{S^{-1}}{R^{-1}} \).
\end{thm}

\begin{proof}
    Let \( R \) and \( S \) be sets of ordered pairs.
    Let \( x_{0} \) be an arbitrary object. Thus, by \cref{dfn:inverse-of-a-set-of-ordered-pairs,dfn:composition-of-two-sets-of-ordered-pairs},
    \begin{align*}
        &x_{0} \in \mleft( \composition{R}{S} \mright)^{-1} \\
        &\iff
        \exists a \, \exists b \, \Big[ x_{0} = \orderedpair{a}{b} \land \orderedpair{b}{a} \in \composition{R}{S} \Big] \\
        &\iff
        \exists a \, \exists b \, \Big[ x_{0} = \orderedpair{a}{b} \land
        \exists c \, \big( \orderedpair{b}{c} \in S \land \orderedpair{c}{a} \in R \big) \Big] \\
        &\iff
        \exists a \, \exists b \, \Big[ x_{0} = \orderedpair{a}{b} \land
        \exists c \, \big( \orderedpair{c}{b} \in S^{-1} \land \orderedpair{a}{c} \in R^{-1} \big)\Big] \\
        &\iff
        \exists a \, \exists b \, \Big[ x_{0} = \orderedpair{a}{b} \land \orderedpair{a}{b} \in \composition{S^{-1}}{R^{-1}} \Big]\\
        &\iff
        x_{0} \in \composition{S^{-1}}{R^{-1}}
    \end{align*}
    Since \( x_{0} \) is arbitrary, the theorem follows by \aref{rule-of-inference:universal-generalization}{Universal Generalization}
    and the \nameref{axm:axiom-of-extensionality}.
\end{proof}

\newpage
\section{Cartesian Products}
\begin{dfn}[Cartesian Product]
    \label{dfn:cartesian-product}
    For any two sets \( A \) and \( B \), the \keyword{cartesian product} of
    \( A \) and \( B \) is the set defined as,
    \begin{equation*}
        A \times B \coloneq \setbuilderbig{\orderedpair{x}{y}}{x \in A \land y \in B}
    \end{equation*}
    More formally,
    \begin{equation*}
        \forall A,B \, \forall x  \, \Big( x \in A \times B \iff
        \exists a,b \, \big[ \, x = \orderedpair{a}{b} \land a \in A \land b \in B \, \big]\Big)
    \end{equation*}
\end{dfn}

\begin{thm}
    \label{thm:existence-of-cartesian-product}
    For any two sets \( A \) and \( B \), \( A \times B \) exists.
\end{thm}

\begin{proof}
    Let \( A \) and \( B \) be arbitrary sets. Let the predicate \( \varphi \) be defined as,
    \begin{equation*}
        \varphi(x,y) \coloneq x \in A \land y \in B
    \end{equation*}
    Let \( a \) and \( b \) be arbitrary objects such that \( \varphi(a,b) \) holds.
    Thus, by the definition of \( \varphi \), it follows that,
    \begin{equation*}
        a \in A \land b \in B
    \end{equation*}
    Clearly, \( a \in A \union B \) and \( b \in A \union B \).
    \begin{equation*}
        \therefore \, \varphi(a,b) \implies a \in A \union B \land b \in A \union B
    \end{equation*}
    Since \( a \) and \( b \) are arbitrary, \aref{rule-of-inference:universal-generalization}{Universal Generalization} yields,
    \begin{equation*}
        \forall x,y \, \Big( \varphi(x,y) \implies x \in A \union B \land y \in A \union B \Big)
    \end{equation*}
    Thus, by \cref{thm:sufficient-condition-for-the-existence-of-sets-of-ordered-pairs-2}, the set
    \( \setbuilder{\orderedpair{x}{y}}{\varphi(x,y)} = \setbuilder{\orderedpair{x}{y}}{x \in A \land y \in B} \)
    exists. By \cref{dfn:cartesian-product}, this set is precisely \( A \times B \).
\end{proof}

\subsection{Properties of Cartesian Products}
Like all other operations involving sets (\( \union, \intersect, -, \symdiff\)), the cartesian product of two sets
obeys certain laws.

\begin{thm}
    \label{thm:non-empty-sets-have-non-empty-cartesian-products}
    For any two sets \( A \) and \( B \), \( A \times B = \emptyset \) if and only if \( A = \emptyset \)
    or \( B = \emptyset \).
    \begin{equation*}
        \forall A,B \, \mleft( A \times B = \emptyset \iff A = \emptyset \lor B = \emptyset \mright)
    \end{equation*}
\end{thm}

\noindent We will give two proofs here.
\begin{proof}[Proof 1.]
    Let \( A \) and \( B \) be arbitrary sets such that \( A \times B = \emptyset \). Suppose it is {\bfseries not}
    the case that \( A = \emptyset \lor B = \emptyset \). In other words, \( A \neq \emptyset \) and
    \( B \neq \emptyset \). By \cref{dfn:empty-set}, this implies,
    \begin{equation*}
        \exists x \, \mleft( x \in A \mright)
        \land
        \exists x \, \mleft( x \in B \mright)
    \end{equation*}
    Let \( a \) and \( b \) be two objects such that \( a \in A \) and \( b \in B \). But by \cref{dfn:cartesian-product},
    this would mean that \( \orderedpair{a}{b} \in A \times B \), which contradicts the fact that \( A \times B = \emptyset \).
    Therefore, by Reduction ad Absurdum, the supposition that \( A \neq \emptyset \land B \neq \emptyset \) must be
    false.
    \begin{equation*}
        \therefore \, A = \emptyset \lor B = \emptyset
    \end{equation*}
    Thus, by the \aref{rule-of-inference:deduction-theorem}{Deduction Theorem}, we derive,
    \begin{equation}
        \label{eq:empty-cartesian-product-implies-at-least-one-empty-set}
        \mleft( A \times B = \emptyset \mright) \implies \mleft( A = \emptyset \lor B = \emptyset \mright)
    \end{equation}
    To prove the converse, suppose that \( A = \emptyset \lor B = \emptyset \).
    Let \( x_{0} = \orderedpair{a}{b} \) be an arbitrary ordered pair.
    If \( A = \emptyset \), then clearly \( a \notin A \). Thus, by \cref{dfn:cartesian-product},
    it follows that \( x_{0} \notin A \times B \). Since \( x_{0} \) is arbitrary,
    we conclude that \( \forall x \, \mleft( x \notin A \times B \mright) \). Thus, \( A \times B = \emptyset \)
    by \cref{dfn:empty-set}. If \( B = \emptyset \), the exact same argument applies. Since \( B = \emptyset \),
    \( b \notin B \) for any \( x_{0} = \orderedpair{a}{b} \). Thus, \( x_{0} \notin A \times B \) for any arbitrary
    \( x_{0} \), so \( A \times B = \emptyset \).
    \begin{equation}
        \label{eq:at-least-one-empty-set-implies-empty-cartesian-product}
        \mleft( A = \emptyset \lor B = \emptyset \mright)
        \implies
        \mleft( A \times B = \emptyset \mright)
    \end{equation}
    Having derived \eqref{eq:empty-cartesian-product-implies-at-least-one-empty-set} and \eqref{eq:at-least-one-empty-set-implies-empty-cartesian-product},
    we can now conclude the biconditional,
    \begin{equation*}
        \mleft( A = \emptyset \lor B = \emptyset \mright)
        \iff
        \mleft( A \times B = \emptyset \mright)
    \end{equation*}
    Since \( A \) and \( B \) are arbitrary, the theorem follows by \aref{rule-of-inference:universal-generalization}{Universal Generalization}.
\end{proof}

\begin{proof}[Proof 2.]
    Let \( A \) and \( B \) be arbitrary sets. By \cref{dfn:cartesian-product,dfn:empty-set},
    \begin{align}
        \label{eq:A-cross-B-equals-empty-predicate-form}
        \mleft( A \times B = \emptyset \mright)
        &\iff
        \forall x \, \big( \lnot \, \exists a,b \, \big[\, x = \orderedpair{a}{b} \land a \in A \land b \in B \,\big] \big) \notag \\
        &\iff
        \forall x \, \forall a,b \, \big( x = \orderedpair{a}{b} \implies \mleft( a \notin A \lor b \notin B \mright) \big)
    \end{align}
    Define the predicates \( \Gamma \) and \( \varphi \) as follows:
    \begin{align}
        \varphi(x,y) &\coloneq x \notin A \lor y \notin B \label{eq:define-inner-predicate} \\
        \Gamma(x)    &\coloneq \forall a,b \, \mleft( x = \orderedpair{a}{b} \implies \varphi(a,b) \mright) \label{eq:define-outer-predicate}
    \end{align}
    \eqref{eq:A-cross-B-equals-empty-predicate-form} can thus be re-expressed as,
    \begin{equation}
        \label{eq:A-cross-B-equals-empty-predicate-expressed-in-terms-of-Gamma}
        \mleft( A \times B = \emptyset \mright) \iff \forall x \, \Gamma(x)
    \end{equation}
    Now, we will use \eqref{eq:define-outer-predicate} to prove that,
    \begin{equation*}
        \forall x \, \Gamma(x) \iff \forall a \, \forall b \, \varphi(a,b)
    \end{equation*}
    Suppose that \( \forall x \, \Gamma(x) \). Let \( \alpha \) and \( \beta \) be arbitrary objects. Since \( \forall x \, \Gamma(x) \),
    then \( \Gamma(\orderedpair{\alpha}{\beta}) \) must hold by \aref{rule-of-inference:universal-instantiation}{Universal Instantiation}. By \eqref{eq:define-outer-predicate},
    this implies,
    \begin{equation*}
        \forall a,b \, \mleft( \orderedpair{\alpha}{\beta} = \orderedpair{a}{b} \implies \varphi(a,b) \mright)
    \end{equation*}
    By \aref{rule-of-inference:universal-instantiation}{Universal Instantiation}, we have,
    \begin{equation*}
        \orderedpair{\alpha}{\beta} = \orderedpair{\alpha}{\beta} \implies \varphi(\alpha, \beta)
    \end{equation*}
    Since \( \orderedpair{\alpha}{\beta} = \orderedpair{\alpha}{\beta} \) is trivially true, then
    \( \varphi(\alpha, \beta) \) must also be true by \aref{rule-of-inference:modus-ponens}{Modus Ponens}. Since \( \alpha \) and \( \beta \)
    are arbitrary, we have,
    \begin{equation*}
        \forall a,b \, \varphi(a,b)
    \end{equation*}
    And by the \aref{rule-of-inference:deduction-theorem}{Deduction Theorem},
    \begin{equation*}
        \forall x \, \Gamma(x) \implies \forall a,b \, \varphi(a,b)
    \end{equation*}
    To prove the converse, suppose that \( \forall a,b \, \varphi(a,b) \) is true. Let \( x_{0} \) be
    an arbitrary object. Obviously, \( x_{0} \) is either an ordered pair or it isn't an ordered pair.
    If \( x_{0} \) is {\bfseries not} an ordered pair, then \( \forall a,b \, \mleft( x_{0} \neq \orderedpair{a}{b} \mright) \).
    As a result, \( \forall a,b \, \mleft( x_{0} = \orderedpair{a}{b} \implies \varphi(a,b) \mright) \) would be
    vacuously  true. Thus, \( \Gamma(x_{0}) \) holds by \eqref{eq:define-outer-predicate}. If \( x_{0} \) is an ordered pair,
    then clearly \( x_{0} = \orderedpair{\alpha}{\beta} \) for some \( \alpha \) and \( \beta \). Since \( \forall a,b \, \varphi(a,b)\),
    it follows that \( \varphi(\alpha,\beta) \) must hold. Let \( a_{0} \) and \( b_{0} \) be arbitrary objects.
    If \( x_{0} = \orderedpair{\alpha}{\beta} = \orderedpair{a_{0}}{b_{0}} \), then \( \alpha = a_{0} \) and \( \beta = b_{0} \) by
    \cref{thm:equality-of-ordered-pairs}. Since \( \varphi(\alpha, \beta) \) is true, it follows that \( \varphi(a_{0}, b_{0}) \) is also true.
    Given that \( a_{0} \) and \( b_{0} \) are arbitrary,
    we have by \aref{rule-of-inference:universal-generalization}{Universal Generalization} \( \forall a,b \, \mleft( x_{0} = \orderedpair{a}{b} \implies \varphi(a,b) \mright) \).
    In other words, \( \Gamma(x_{0}) \) holds.

    \vspace{0.5cm}
    \noindent
    We have shown that \( \Gamma(x_{0}) \) follows from \( \forall a,b \, \varphi(a,b) \) regardless
    of whether \( x_{0} \) is an ordered pair. Since \( x_{0} \) is arbitrary, we can conclude by \aref{rule-of-inference:universal-generalization}{Universal Generalization} that
    \( \forall x \, \Gamma(x) \). Therefore,
    \begin{equation*}
        \forall a,b \, \varphi(a,b) \implies \forall x \, \Gamma(x)
    \end{equation*}
    Having proven both implications, we are now in the position to conclude that,
    \begin{equation*}
        \forall x \, \Gamma(x) \iff \forall a,b \, \varphi(a,b)
    \end{equation*}
    From \eqref{eq:A-cross-B-equals-empty-predicate-expressed-in-terms-of-Gamma}, we thus have,
    \begin{equation}
        \label{eq:A-cross-B-equals-empty-final-statement}
        \mleft( A \times B = \emptyset \mright) \iff \forall x,y \, \varphi(a,b) \iff \forall x,y \, \mleft( x \notin A \lor y \notin B \mright)
    \end{equation}
    Using \cref{tautology:grouping-disjuncts-into-universally-quantified-formula}, we can simplify \eqref{eq:A-cross-B-equals-empty-final-statement}
    even further.
    \begin{align*}
        \mleft( A \times B = \emptyset \mright)
        &\iff
        \forall x \, \mleft[\, \forall y \, \mleft( x \notin A \lor y \notin B \mright) \,\mright] \\
        &\iff
        \forall x \, \mleft[\, x \notin A \lor \forall y \, \mleft( y \notin B \mright) \,\mright]\\
        &\iff
        \forall x \, \big( x \notin A \big) \lor \forall y \, \mleft( y \notin B \mright) \\
        &\iff
        \forall x \, \mleft( x \notin A \mright) \lor \forall x \, \mleft( x \notin B \mright) \\
        &\iff
        A = \emptyset \lor B = \emptyset
    \end{align*}
    Since \( A \) and \( B \) are arbitrary, the theorem follows by \aref{rule-of-inference:universal-generalization}{Universal Generalization}.
\end{proof}

\begin{thm}
    \label{thm:relationship-between-subset-and-cartesian-products}
    If \( S \) is a non-empty set, then \( A \subseteq B \) if and only if \( \mleft( A \times S \mright) \subseteq \mleft( B \times S \mright) \) for any \( A \) and \( B \).
    \begin{equation*}
        \forall S \, \Big( S \neq \emptyset \implies \forall A,B \, \big( A \subseteq B \iff \mleft( A \times S \mright) \subseteq \mleft( B \times S \mright)\big) \Big)
    \end{equation*}
\end{thm}

\begin{proof}
    Let \( A \), \( B \), and \( S \) be arbitrary sets such that \( S \neq \emptyset \). Suppose that \( A \subseteq B \). By \cref{dfn:subset},
    we have,
    \begin{equation}
        \label{eq:every-element-of-A-is-also-an-element-of-B}
        \forall x \, \mleft( x \in A \implies x \in B \mright)
    \end{equation}
    Let \( x_{0} \in \mleft( A \times S \mright) \) be arbitrary.
    Thus \( x_{0} = \orderedpair{a_{0}}{b_{0}} \)
    for some \( a_{0} \in A \) and \( b_{0} \in S \) by \cref{dfn:cartesian-product}.
    By \eqref{eq:every-element-of-A-is-also-an-element-of-B}, necessarily \( a_{0} \in B \).
    Thus, \( x_{0} = \orderedpair{a_{0}}{b_{0}} \in \mleft( B \times S \mright) \).
    \begin{equation*}
        \therefore \, x_{0} \in \mleft( A \times S \mright) \implies x_{0} \in \mleft( B \times S \mright)
    \end{equation*}
    Since \( x_{0} \) is arbitrary, it follows that,
    \begin{equation*}
        \forall x \, \Big( x \in \mleft( A \times S \mright) \implies x \in \mleft( B \times S \mright) \Big)
    \end{equation*}
    Thus \( \mleft( A \times S \mright) \subseteq \mleft( B \times S \mright) \) by \cref{dfn:subset}.
    \begin{equation}
        \label{eq:A-subset-B-implies-A-times-S-subset-B-times-S}
        \therefore \,
        A \subseteq B \implies \mleft( A \times S \mright) \subseteq \mleft( B \times S \mright)
    \end{equation}
    To prove the converse, suppose that \( \mleft( A \times S \mright) \subseteq \mleft( B \times S \mright) \).
    Let \( x_{0} \in A \) be arbitrary. Since \( S \neq \emptyset \), it follows that there exists a \( z \in S \).
    Since \( x_{0} \in A \land z \in S \), it follows that \( \orderedpair{x_{0}}{z} \in \mleft( A \times S \mright) \).
    And since \( \mleft( A \times S \mright) \subseteq \mleft( B \times S \mright) \),
    \( \orderedpair{x_{0}}{z} \in \mleft( B \times S \mright) \) must also be true. By \cref{dfn:cartesian-product},
    \( \orderedpair{x_{0}}{z} \in \mleft( B \times S \mright) \iff x_{0} \in B \land z \in S \).
    Thus, we infer \( x_{0} \in B \) by \aref{rule-of-inference:modus-ponens}{Modus Ponens}. By the \aref{rule-of-inference:deduction-theorem}{Deduction Theorem},
    \begin{equation*}
        x_{0} \in A \implies x_{0} \in B
    \end{equation*}
    Since \( x_{0} \) is arbitrary, we conclude by \aref{rule-of-inference:universal-generalization}{Universal Generalization} that,
    \begin{equation*}
        \forall x \, \mleft( x \in A \implies x \in B \mright)
    \end{equation*}
    By \cref{dfn:subset}, it follows that \( A \subseteq B\).
    \begin{equation}
        \label{eq:A-times-S-subset-B-times-S-implies-A-subset-B}
        \therefore \, \mleft( A \times S \mright) \subseteq \mleft( B \times S \mright) \implies A \subseteq B
    \end{equation}
    And from \eqref{eq:A-subset-B-implies-A-times-S-subset-B-times-S} and \eqref{eq:A-times-S-subset-B-times-S-implies-A-subset-B},
    we infer,
    \begin{equation*}
        A \subseteq B \iff \mleft( A \times S \mright) \subseteq \mleft( B \times S \mright)
    \end{equation*}
    Thus, by the \aref{rule-of-inference:deduction-theorem}{Deduction Theorem},
    \begin{equation*}
        S \neq \emptyset \implies \Big( A \subseteq B \iff \mleft( A \times S \mright) \subseteq \mleft( B \times S \mright) \Big)
    \end{equation*}
    Since \( A \), \( B \), and \( S \) are arbitrary, the theorem follows from \aref{rule-of-inference:universal-generalization}{Universal Generalization}.
\end{proof}

\begin{thm}[Right Distributivity of \( \times \) over \( \union, \intersect, -, \symdiff \)]
    \label{set-identities:right-distributivity-of-cartesian-products}
    For any three sets \( A \), \( B \) and \( C \), the following identities hold:
    \begin{gather*}
        \mleft( A \union B \mright) \times C = \mleft( A \times C \mright) \union \mleft( B \times C \mright)\\
        \mleft( A \intersect B \mright) \times C = \mleft( A \times C \mright) \intersect \mleft( B \times C \mright)\\
        \mleft( A - B \mright) \times C = \mleft( A \times C \mright) - \mleft( B \times C \mright)\\
        \mleft( A \symdiff B \mright) \times C = \mleft( A \times C \mright) \symdiff \mleft( B \times C \mright)
    \end{gather*}
\end{thm}

\begin{proof}
    Let \( A \), \( B \), and \( C \) be arbitrary sets. To prove the first three equalities,
    we shall write down the predicates for the set on each side of the equality, and then
    prove that they are equivalent.

    \vspace{0.3cm}
    \noindent
    {\scshape Claim}: \( \mleft( A \union B \mright) \times C = \mleft( A \times C \mright) \union \mleft( B \times C \mright) \)

    \vspace{0.1cm}
    \noindent
    Start by defining the predicates \( \Gamma \) and \( \varphi \) as follows:
    \begin{align*}
        \Gamma(x)
        &\coloneq
        \exists a \, \exists b \, \big[ x = \orderedpair{a}{b} \land \mleft( a \in A \lor a \in B \mright) \land b \in C \big] \\
        \varphi(x)
        &\coloneq
        \exists a \, \exists b \, \big[ x = \orderedpair{a}{b} \land a \in A \land b \in C \big]
        \lor
        \exists a \, \exists b \, \big[ x = \orderedpair{a}{b} \land a \in B \land b \in C \big]
    \end{align*}
    Thus, for any arbitrary \( x \),
    \begin{equation*}
        \begin{split}
            \Gamma(x)
            &\iff \exists a \, \exists b \, \big[ x = \orderedpair{a}{b} \land \mleft( a \in A \lor a \in B \mright) \land b \in C \big] \\
            &\iff \exists a \, \exists b \, \big[ x = \orderedpair{a}{b} \land
            \big( \mleft( a \in A \land b \in C \mright) \lor \mleft( a \in B \land b \in C \mright) \big) \big] \\
            &\iff \exists a \, \exists b \, \big[ x = \orderedpair{a}{b} \land a \in A \land b \in C \big]
            \lor
            \exists a \, \exists b \, \big[ x = \orderedpair{a}{b} \land a \in B \land b \in C \big]\\
            &\iff \varphi(x)
        \end{split}
    \end{equation*}
    \aref{rule-of-inference:universal-generalization}{Univesal Generalization} then yields \( \forall x \, \mleft( \Gamma(x) \iff \varphi(x) \mright) \).
    However, the following also holds by \cref{thm:union-in-terms-of-logical-or,dfn:set-builder-notation-ordered-pair,dfn:cartesian-product}:
    \begin{gather*}
        \forall x \, \big[ x \in \mleft( A \union B \mright) \times C \iff x \in \Gamma(x) \big]\\
        \forall x \, \big[ x \in \mleft( A \times C \mright) \union \mleft( B \times C \mright) \iff x \in \varphi(x) \big]
    \end{gather*}
    The first identity follows immediately from the \nameref{axm:axiom-of-extensionality}.

    \vspace{0.5cm}
    \noindent
    {\scshape Claim}: \( \mleft( A \intersect B \mright) \times C = \mleft( A \times C \mright) \intersect \mleft( B \times C \mright) \)

    \vspace{0.1cm}
    \noindent
    Define \( \Gamma \) and \( \varphi \) as follows:
    \begin{align*}
        \Gamma(x)
        &\coloneq
        \exists a \, \exists b \, \big[ x = \orderedpair{a}{b} \land \mleft( a \in A \land a \in B \mright) \land b \in C \big] \\
        \varphi(x)
        &\coloneq
        \exists a \, \exists b \, \big[ x = \orderedpair{a}{b} \land \mleft( a \in A \land b \in C \mright) \big]
        \land
        \exists a \, \exists b \, \big[ x = \orderedpair{a}{b} \land \mleft( a \in B \land b \in C \mright) \big]
    \end{align*}
   
    \noindent
    Let \( x_{0} \) be an arbitrary object such that \( \Gamma(x_{0}) \) holds true. It follows from the definition of \( \Gamma \)
    that \( x_{0} = \orderedpair{\alpha}{\beta} \) for some \( \alpha \) and \( \beta \) where \( \alpha \in A \land \alpha \in B \)
    and \( \beta \in C \). Clearly, \( \varphi(x_{0}) \) also holds. Therefore, \( \Gamma(x_{0}) {\implies} \varphi(x_{0}) \). To prove
    the converse, suppose that \( \varphi(x_{0}) \) is true. By the definition of \( \varphi \),
    it follows that \( x_{0} = \orderedpair{\alpha}{\beta} = \orderedpair{\gamma}{\delta}\) where \( \alpha \in A \), \( \beta \in C \),
    \( \gamma \in B \), and \( \delta \in C \). Since \( \orderedpair{\alpha}{\beta} = \orderedpair{\gamma}{\delta} \), necessarily
    \( \alpha = \gamma \) and \( \beta = \delta \) by \cref{thm:equality-of-ordered-pairs}. Since \( \alpha = \gamma \) and
    \( \alpha \in A \land \gamma \in B \), it follows that \( \alpha \in A \land \alpha \in B \) must be true. Therefore,
    \( x_{0} = \orderedpair{\alpha}{\beta} \land (\alpha \in A \land \alpha \in B) \land \beta \in C \) is true. In other words,
    \( \Gamma(x_{0}) \) holds. Thus, \( \varphi(x_{0}) {\implies} \Gamma(x_{0}) \).
    
    \vspace{0.2cm}
    \noindent
    Having established the biconditional \( \Gamma(x_{0}) {\iff} \varphi(x_{0}) \), we conclude by \aref{rule-of-inference:universal-generalization}{Universal Generalization} that,
    \( \forall x \, \mleft( \Gamma(x) {\iff} \varphi(x) \mright) \).
    Furthermore, we have by \cref{thm:intersection-in-terms-of-logical-and,dfn:cartesian-product},
    \begin{gather*}
        \forall x \, \big[ x \in \mleft( A \intersect B \mright) \times C \iff \Gamma(x) \big] \\
        \forall x \, \big[ x \in \mleft( A \times C \mright) \intersect \mleft( B \times C \mright) \iff \varphi(x) \big]
    \end{gather*}
    The second equality follows readily from the \nameref{axm:axiom-of-extensionality}.

    \vspace{0.5cm}
    \noindent
    {\scshape Claim}: \( \mleft( A - B \mright) \times C = \mleft( A \times C \mright) - \mleft( B \times C \mright) \)

    \vspace{0.1cm}
    \noindent
    Define \( \Gamma \) and \( \varphi \) as follows:
    \begin{align*}
        \Gamma(x)
        &\coloneq
        \exists a \, \exists b \, \big[ x = \orderedpair{a}{b} \land \mleft( a \in A \land a \notin B \mright) \land b \in C \big] \\
        \varphi(x)
        &\coloneq
        \exists a \, \exists b \, \big[ x = \orderedpair{a}{b} \land a \in A \land b \in C \big]
        \land
        \forall a \, \forall b \, \big[ x = \orderedpair{a}{b} \implies a \notin B \lor b \notin C \big]
    \end{align*}
    It is not difficult to see that \( \forall x \, \big( x \in \mleft( A - B \mright) \times C {\iff} \Gamma(x) \big) \). The relationship between
    \( \varphi(x) \) and \( \mleft( A \times C \mright) - \mleft( B \times C \mright) \) is more subtle but should be perfectly intelligible after
    an application of the laws of logic. As in, for any arbitrary \( x \),
    \begin{equation*}
        \begin{split}
            &x \in \mleft( A \times C \mright) - \mleft( B \times C \mright)\\
            &\iff \exists a \, \exists b \, \big[ x = \orderedpair{a}{b} \land a \in A \land b \in C \big]
            \land \lnot \,
            \exists a \, \exists b \, \big[ x = \orderedpair{a}{b} \land a \in B \land b \in C \big] \\
            &\iff
            \exists a \, \exists b \, \big[ x = \orderedpair{a}{b} \land a \in A \land b \in C \big]
            \land
            \forall a \, \forall b \, \big[ x \neq \orderedpair{a}{b} \lor \mleft( a \notin B \lor b \notin C \mright) \big]\\
            &\iff
            \exists a \, \exists b \, \big[ x = \orderedpair{a}{b} \land a \in A \land b \in C \big]
            \land
            \forall a \, \forall b \, \big[ x = \orderedpair{a}{b} \implies \mleft( a \notin B \lor b \notin C \mright) \big]\\
            &\iff
            \varphi(x)
        \end{split}
    \end{equation*}
    Thus, it remains to prove that \( \forall x \, \mleft( \Gamma(x) {\iff} \varphi(x) \mright) \). Let \( x_{0} \) be an arbitrary object
    such that \( \Gamma(x_{0}) \) is satisfied. Thus, \( x_{0} = \orderedpair{\alpha}{\beta} \) for some \( \alpha \) and \( \beta \)
    such that \( \alpha \in A \land \alpha \notin B \) and \( \beta \in C \). Since \( \alpha \in A \) and \( \beta \in C \) and
    \( x_{0} = \orderedpair{\alpha}{\beta}\), it follows that,
    \begin{equation*}
        \exists a \, \exists b \, \big[ x_{0} = \orderedpair{a}{b} \land a \in A \land b \in C \big]
    \end{equation*}
    Let \( a_{0} \) and \( b_{0} \) be arbitrary objects such that \( x_{0} = \orderedpair{a_{0}}{b_{0}} \). Since \( x_{0} = \orderedpair{\alpha}{\beta} \),
    it follows from \cref{thm:equality-of-ordered-pairs} that \( a_{0} = \alpha \) and \( b_{0} = \beta \). Since \( a_{0} = \alpha \) and
    \( \alpha \notin B \), it follows that \( a_{0} \notin B \).
    By \aref{rule-of-inference:disjunction-introduction}{Disjunction Introduction}, we conclude \( a_{0} \notin B \lor b_{0} \notin C \).
    In other words, \( (x_{0} = \orderedpair{a_{0}}{b_{0}}) {\implies} \mleft( a_{0} \notin B \lor b_{0} \notin C \mright) \) for any arbitrary
    \( a_{0} \) and \( b_{0} \). By \aref{rule-of-inference:universal-generalization}{Universal Generalization}, we have,
    \begin{equation*}
        \forall a \, \forall b \, \big[ x = \orderedpair{a}{b} \implies \mleft( a \notin B \lor b \notin C \mright) \big]
    \end{equation*}
    By \aref{rule-of-inference:conjunction-introduction}{Conjunction Introduction}, it follows that,
    \begin{equation*}
        \exists a \, \exists b \, \big[ x_{0} = \orderedpair{a}{b} \land a \in A \land b \in C \big]
        \land
        \forall a \, \forall b \, \big[ x_{0} = \orderedpair{a}{b} {\implies} \mleft( a \notin B \lor b \notin C \mright) \big]
    \end{equation*}
    In other words, we have just derived \( \varphi(x_{0}) \) from the supposition of \( \Gamma(x_{0}) \)'s truth.
    \begin{equation*}
        \therefore \, \Gamma(x_{0}) \implies \varphi(x_{0})
    \end{equation*}

    \vspace{0.2cm}
    \noindent
    To prove the converse, suppose that \( \varphi(x_{0}) \) holds. By the definition of \( \varphi \),
    \begin{equation*}
        \exists a \, \exists b \, \big[ x_{0} = \orderedpair{a}{b} \land a \in A \land b \in C \big]
        \land
        \forall a \, \forall b \, \big[ x_{0} = \orderedpair{a}{b} {\implies} \mleft( a \notin B \lor b \notin C \mright) \big]
    \end{equation*}
    The left conjunct implies that \( x_{0} = \orderedpair{\alpha}{\beta}\) for some \( \alpha \) and \( \beta \)
    where \( \alpha \in A \) and \( \beta \in C \). On the other hand, the right conjunct implies,
    \begin{equation*}
        \alpha \notin B \lor \beta \notin C
    \end{equation*}
    But it has already been established that \( \beta \in C \), so by \aref{rule-of-inference:disjunctive-syllogism}{Disjunctive Syllogism}, \( \alpha \notin B \) must be
    true. Therefore, \( x_{0} = \orderedpair{\alpha}{\beta} \land \alpha \in A \land \alpha \notin B \land \beta \in C \),
    so \( \Gamma(x_{0}) \) is indeed satisfied.
    \begin{equation*}
        \therefore \, \varphi(x_{0}) \implies \Gamma(x_{0})
    \end{equation*}
    Since the implication holds in both directions, we derive via \aref{rule-of-inference:biconditional-introduction}{Biconditional Introduction}
    \( \Gamma(x_{0}) {\iff} \varphi(x_{0}) \) and subsequently
    \( \forall x \, \mleft( \Gamma(x) {\iff} \varphi(x) \mright) \) by \aref{rule-of-inference:universal-generalization}{Universal Generalization}.
    The third equality follows readily from the \nameref{axm:axiom-of-extensionality}.

    \vspace{0.5cm}
    \noindent
    {\scshape Claim}: \( \mleft( A \symdiff B \mright) \times C = \mleft( A \times C \mright) \symdiff \mleft( B \times C \mright) \)

    \vspace{0.1cm}
    \noindent
    This equality can be easily proven using the first and third equalities above.
    \begin{align*}
        \mleft( A \symdiff B \mright) \times C
        &=
        \big[ \mleft( A - B \mright) \union \mleft( B - A \mright) \big] \times C \\
        &=
        \big[  \mleft( A - B \mright) \times C \big] \union \big[ \mleft( B - A \mright) \times C \big] \\
        &=
        \big[ \mleft( A \times C \mright) - \mleft( B \times C \mright) \big]
        \union
        \big[ \mleft( B \times C \mright) - \mleft( A \times C \mright) \big] \\
        &=
        \mleft( A \times C \mright) \symdiff \mleft( B \times C \mright) \qedhere
    \end{align*}
\end{proof}

\begin{rem}
    The \keyword{left-distributivity} of cartesian products can be proven in the exact same way as above.
\end{rem}

\newpage
\section{Set-theoretic Definition of Functions}
\begin{dfn}
    \label{dfn:function-basic}
    A set \( f \) is a \keyword{function} if and only if the following statement holds:
    \begin{equation*}
        \forall x \, \Big( x \in f \implies \exists a \, \exists b \, \mleft( x = \orderedpair{a}{b} \mright) \Big)
        \land
        \forall x \, \forall y \, \forall z \, \Big( {\orderedpair{x}{y}} \in f \land \orderedpair{x}{z} \in f \implies y = z \Big)
    \end{equation*}
\end{dfn}

\begin{rem}
    A function is a set of ordered pairs such that the first coordinate of every ordered pair it contains is unique.
    It isn't difficult to see that this definition is consistent with how we intuitively think about functions---as rules
    that associate an object with a single unique object. Also, since functions are sets of ordered pairs, all the properties
    presented in \cref{ssec:properties-of-sets-of-ordered-pairs} also apply to functions.
\end{rem}


\begin{dfn}[Arrow Notation]
    \label{dfn:arrow-notation}
    Let \( A \), \( B \), and \( f \) be  sets. We use the notation
    \function{f}{A}{B} if and only if \( f \) is a function with \( \dom f = A \) and \( \ran f \subseteq B \).
    Informally, the notation can be read as ``\( f \) is a function on \( A \) into \( B \).''
\end{dfn}

\begin{dfn}[Functional Notation]
    \phantomsection \label{dfn:functional-notation}
    For any function \( f \), if \( \orderedpair{x}{y} \in f \) for some \( x \in \dom f \), then
    \( y \) is called the \keyword{value} of \( f \) at \( x \) and is denoted by \( f(x) \).
    This can be stated in first-order logic as,
    \begin{equation*}
        \forall x \, \forall y \, \Big(
        \orderedpair{x}{y} \in f
        \iff
        y = f(x) \Big)
    \end{equation*}
\end{dfn}

\begin{cor}
    \label{thm:express-the-range-of-a-function-using-functional-notation}
    For any function \( f \), \( \ran f = \setbuilderbig{y}{\exists x \, \mleft( y = f(x) \mright)} \)
\end{cor}

\begin{dfn}[Index Notation]
    \label{dfn:index-notation}
    The value of a function \( f \) at \( x \) can also be denoted as \( f_{x} \).
\end{dfn}

\begin{dfn}[Injection]
    \label{dfn:injection}
    A function \( f \) is \keyword{injective} if and only if,
    \begin{equation*}
        \forall x \, \forall y \, \Big( x \in \dom f \land y \in \dom f \implies \big( f(x) = f(y) \implies x = y \big)\Big)
    \end{equation*}
    We use the notation \injection{f}{A}{B} to mean,
    ``\( f \) is a injective/one-to-one function on \( A \) into \( B \),''
    or
    ``\( f \) is an injection from \( A \) to \( B \).''
\end{dfn}

\begin{dfn}[Surjection]
    \label{dfn:surjection}
    A function \function{f}{A}{B} is \keyword{surjective} if and only if,
    \begin{equation*}
        \forall y \, \mleft( y \in B \iff \exists x \, \mleft( x \in A \land \orderedpair{x}{y} \in f \mright) \mright)
    \end{equation*}
    In other words, \function{f}{A}{B} is surjective if and only if \( \ran f = B \).
    The notation \surjection{f}{A}{B} reads,
    ``\( f \) is a surjective function on \( A \) into \( B \),''
    ``\( f \) is a function on \( A \) onto \( B \),'' or
    ``\( f \) is a surjection from \( A \) to \( B \).''
\end{dfn}

\begin{cor}
    \label{thm:every-function-on-its-domain-into-its-range-is-surjective}
    If \( f \) is a function, then \function{f}{\dom f}{\ran f} is surjective.
    \begin{equation*}
        \forall A,B \, \forall f \, \Big( \function{f}{A}{B} \implies \surjection{f}{A}{\ran f} \Big)
    \end{equation*}
\end{cor}

\begin{rem}
    The notion of ``surjection'' is only relevant in contexts where a function is described in relation
    to another set---typically a set that contains the function's range. Incidentally, this is also the
    sole purpose for the usage of the \nameref{dfn:arrow-notation}. A function can be fully described
    without the use of the \nameref{dfn:arrow-notation}; after all, a function is simply a set
    of ordered pairs. The \nameref{dfn:arrow-notation} is only useful when one intends to describe the
    relationship between the function and another set.
\end{rem}

\begin{dfn}[Bijection]
    \label{dfn:bijection}
    A function \( f \) is \keyword{bijective} if and only if
    it is both injective and surjective.
    We use the notation \bijection{f}{A}{B} to mean, ``\( f \) is a one-to-one
    function on \( A \) onto \( B \),''
    or ``\( f \) is a bijection between \( A \) and \( B \).''
\end{dfn}

\begin{cor}
    \label{thm:every-injective-function-on-its-domain-into-its-range-is-a-bijection}
    If \( f \) is an injective function, then \function{f}{\dom f}{\ran f} is bijective.
    \begin{equation*}
        \forall A,B \, \forall f \, \Big( \injection{f}{A}{B} \implies \bijection{f}{A}{\ran f} \Big)
    \end{equation*}
\end{cor}

\begin{dfn}[Permutation]
    \phantomsection \label{dfn:permutation-of-a-set}
    Given any set \( A \), any function \( \bijection{\rho}{A}{A} \) is called a \keyword{permutation} of \( A \).
\end{dfn}

\begin{dfn}[Invertibility]
    \label{dfn:invertibility-of-functions}
    A function \( f \) is \keyword{invertible} if and only if \( f^{-1} \) is also a function.
\end{dfn}

\begin{rem}
    Although \( f^{-1} \) exists for any function \( f \)
    (\cref{thm:existence-of-the-inverse-of-a-set-of-ordered-pairs}),
    it is not necessarily a function. Unless it has been proven that
    \( f^{-1} \) satisfies \cref{dfn:function-basic}, \( f^{-1} \) should
    be treated as nothing more than a set of ordered pairs.
\end{rem}

\begin{dfn}[Restriction]
    \label{dfn:restricted-functions}
    If \( f \) is a function and \( A \) is a set, then the \keyword{restriction} of \( f \) to \( A \)
    is the set given by,
    \begin{equation*}
        \restrictedfunction{f}{A} \coloneq \setbuilderbig{\orderedpair{x}{y} \in f}{x \in A}
    \end{equation*}
    We also say that \( f \) is an \keyword{extension} of \restrictedfunction{f}{A}.
\end{dfn}

\begin{dfn}[Partial Function]
    \label{dfn:partial-function}
    A function \( f \) is a \keyword{partial function} from \( A \) to \( B \) if and only if
    there exists some \( D \subseteq A \) such that \( f \) is a function on \( D \) into \( B \).
    The notation \( \partialfunction{f}{A}{B} \) reads, ``\( f \) is a partial function from \( A \) to \( B \).''
    \begin{equation*}
        \forall f \, \forall A \, \forall B \,
        \Big[ \partialfunction{f}{A}{B} \iff \exists D \, \big( D \subseteq A \land \function{f}{D}{B} \big) \Big]
    \end{equation*}
\end{dfn}

\begin{rem}
    Obviously, \( \function{f}{A}{B} \implies \partialfunction{f}{A}{B} \), though the converse isn't true.
\end{rem}

\vspace{0.2cm}
\noindent
Partial functions can also be surjections, injections, or bijections, which leads us to the following definitions.

\begin{dfn}[Partial Injection]
    \label{dfn:partial-injection}
    A function \( f \) is a \keyword{partial injection} from \( A \) to \( B \) if and only if there exists some \( D \subseteq A \)
    such that \( f \) is an injection from \( D \) to \( B \). The notation \( \partialinjection{f}{A}{B} \) reads,
    ``\( f \) is a partial injection from \( A \) to \( B \).''
    \begin{equation*}
        \forall f \, \forall A \, \forall B \, \Big[ \partialinjection{f}{A}{B} \iff \exists D \, \big( D \subseteq A \land \injection{f}{D}{B} \big) \Big]
    \end{equation*}
\end{dfn}

\begin{dfn}[Partial Surjection]
    \label{dfn:partial-surjection}
    A function \( f \) is a \keyword{partial surjection} from \( A \) to \( B \) if and only if there exists some \( D \subseteq A \)
    such that \( f \) is a surjection from \( D \) to \( B \). The notation \( \partialsurjection{f}{A}{B} \) reads,
    ``\( f \) is a partial surjection from \( A \) to \( B \).''
    \begin{equation*}
        \forall f \, \forall A \, \forall B \, \Big[ \partialsurjection{f}{A}{B} \iff \exists D \, \big( D \subseteq A \land \surjection{f}{D}{B} \big) \Big]
    \end{equation*}
\end{dfn}

\begin{dfn}[Partial Bijection]
    \phantomsection \label{dfn:partial-bijection}
    A function \( f \) is a \keyword{partial bijection} from \( A \) to \( B \)
    if and only if there exists some \( D \subseteq A \)
    such that \( f \) is a bijection between \( D \) and \( B \).
    The notation \( \partialbijection{f}{A}{B} \) reads,
    ``\( f \) is a partial bijection from \( A \) to \( B \).''
    \begin{equation*}
        \forall f \, \forall A \, \forall B \, \Big[ \partialbijection{f}{A}{B} \iff \exists D \, \big( D \subseteq A \land \bijection{f}{D}{B} \big) \Big]
    \end{equation*}
\end{dfn}

\begin{dfn}[Identity Function]
    \phantomsection \label{dfn:identity-function}
    For any set \( A \), the identity function on \( A \) is defined as follows:
    \begin{equation*}
        \identityfunc{A} = \setbuilderBig{\orderedpair{x}{y}}{ x \in A \,\,{\land}\,\, y \in A \,\,{\land}\,\, x = y}
    \end{equation*}
\end{dfn}

\begin{cor}
    \phantomsection \label{cor:identity-function-of-empty-set-is-the-empty-set}
    \( \identityfunc{\emptyset} = \emptyset \).
\end{cor}

\begin{cor}
    \phantomsection \label{cor:identity-function-is-a-bijection}
    For any set \( A \), \( \bijection{\identityfunc{A}}{A}{A} \).
\end{cor}

\begin{dfn}[System of Functions]
    \label{dfn:system-of-functions}
    A non-empty set \( F \) is called a \keyword{system} or \keyword{family} of functions
    if and only if every element of \( F \) is a function.
\end{dfn}

\begin{dfn}[Compatibility]
    \label{dfn:compatibility-of-a-system-of-functions}
    A system of functions \( F \) is a \keyword{compatible} if and only if,
    \begin{equation*}
        \forall x \,
        \Big( x \in \Intersect\setbuilderBig{\dom f}{f \in F}
        \implies
        \forall f \, \forall g \, \big( f \in F \land g \in F \implies f(x) = g(x) \big) \Big)
    \end{equation*}
\end{dfn}

\begin{cor}
    \label{thm:compatibility-of-two-functions}
    Two functions \( f \) and \( g \) are compatible if and only if,
    \begin{equation*}
        \forall x \, \Big[ x \in \mleft( \dom f \intersect \dom g \mright) \implies f(x) = g(x) \Big]
    \end{equation*}
\end{cor}

\subsection{Properties of Functions}
\begin{thm}
    \label{thm:a-function-is-a-subset-of-the-cartesian-product-between-its-domain-and-range}
    If \( f \) is a function, then \( f \subseteq \mleft( \dom f \times \ran f \mright) \)
\end{thm}

\begin{proof}
    Let \( f \) be a function. By \cref{dfn:cartesian-product},
    \begin{equation}
        \label{eq:cartesian-product-of-domain-and-range}
        \dom f \times \ran f = \setbuilderbig{\orderedpair{x}{y}}{x \in \dom f \land y \in \ran f}
    \end{equation}
    Suppose that \( \alpha \in f \). Since \( f \) is a function, it follows from \cref{dfn:function-basic}
    that,
    \begin{equation*}
        \begin{aligned}
            \alpha = \orderedpair{x_{0}}{y_{0}} &
        \end{aligned}
        \begin{aligned}
            &\text{for some \( x_{0} \) and \( y_{0} \)}
        \end{aligned}
    \end{equation*}
    By \cref{dfn:domain-of-a-set-of-ordered-pairs,dfn:range-of-a-set-of-ordered-pairs}, clearly
    \( x_{0} \in \dom f\) and \( y_{0} \in \ran f \). By \eqref{eq:cartesian-product-of-domain-and-range},
    it follows that \( \alpha = \orderedpair{x_{0}}{y_{0}} \in \mleft( \dom f \times \ran f \mright) \).
    Since \( \alpha \) is arbitrary, we conclude by \aref{rule-of-inference:universal-generalization}{Universal Generalization} that,
    \begin{equation*}
        \forall x \, \big( x \in f \implies x \in \mleft( \dom f \times \ran f \mright) \big)
    \end{equation*}
    By \cref{dfn:subset}, this means that \( f \subseteq \mleft( \dom f \times \ran f \mright) \).
\end{proof}

\begin{cor}
    \label{thm:a-function-on-A-into-B-is-a-subset-of-the-cartesian-product-between-A-and-B}
    If \function{f}{A}{B}, then \( f \subseteq A \times B \) because \( \dom f = A \) and \( \ran f \subseteq B \)
    so that \( f \subseteq \mleft( A \times \ran f \mright) \subseteq \mleft( A \times B \mright) \).
\end{cor}

\begin{thm}
    \label{thm:full-definition-of-the-arrow-notation}
    For any three sets \( f \), \( A \), and \( B \),
    \begin{multline*}
        \function{f}{A}{B}
        \iff
        \bigg(
        \forall x \, \Big[ x \in f \implies \exists a,b \, \big( x = \orderedpair{a}{b} \land a \in A \land b \in B \big) \Big] \\
        \land
        \forall x \, \Big[ x \in A \implies \exists! y \, \big( y \in B \land \orderedpair{x}{y} \in f \big) \Big]\bigg)
    \end{multline*}
\end{thm}

\begin{proof}
    Let \( f \), \( A \), and \( B \) be arbitrary sets. Suppose that \function{f}{A}{B}. That is, \( f \) is a function with
    \( \dom f = A \) and \( \ran f \subseteq B \). Let \( x_{0} \) be an arbitrary object such that \( x_{0} \in f \). By
    \cref{dfn:function-basic}, it follows that,
    \begin{equation*}
        \begin{aligned}
            x_{0} = \orderedpair{\alpha}{\beta}&
        \end{aligned}
        \begin{aligned}
            &\text{for some \( \alpha \) and \( \beta \)}
        \end{aligned}
    \end{equation*}
    Since \( x_{0} = \orderedpair{\alpha}{\beta} \) and \( x_{0} \in f \), it follows from \cref{dfn:domain-of-a-set-of-ordered-pairs,dfn:range-of-a-set-of-ordered-pairs}
    that \( \alpha \in \dom f \) and \( \beta \in \ran f \). And since \( \dom f = A \) and \( \ran f \subseteq B\), it follows that \( \alpha \in A \) and \( \beta \in B \).
    \begin{equation}
        \label{eq:first-conjunct}
        \therefore \, \forall x \, \Big[ x \in f \implies \exists a,b \, \big( x = \orderedpair{a}{b} \land a \in A \land b \in B \big) \Big]
    \end{equation}
    Now, let \( x_{0} \in A \) be arbitrary. Since \( \dom f = A \), \( x_{0} \in \dom f \) by the \nameref{axm:axiom-of-extensionality}.
    By \cref{dfn:domain-of-a-set-of-ordered-pairs,dfn:range-of-a-set-of-ordered-pairs}, this means that,
    \begin{equation*}
        \begin{aligned}
            \orderedpair{x_{0}}{y_{0}} \in f&
        \end{aligned}
        \begin{aligned}
            &\text{for some \( y_{0} \in \ran f \)}
        \end{aligned}
    \end{equation*}
    And since \( \ran f \subseteq B \), it follows from \cref{dfn:subset} that \( y_{0} \in B \). Suppose there exists another \( z_{0} \) such that
    \( \orderedpair{x_{0}}{z_{0}} \in f \). But \( f \) is a function, so,
    \begin{equation*}
        \orderedpair{x_{0}}{y_{0}} \in f \land \orderedpair{x_{0}}{z_{0}} \in f \implies y_{0} = z_{0}
    \end{equation*}
    by \cref{dfn:function-basic}. Thus, \( y_{0} \) is unique.
    By the \aref{rule-of-inference:deduction-theorem}{Deduction Theorem} and \aref{rule-of-inference:universal-generalization}{Universal Generalization}, we infer,
    \begin{equation}
        \label{eq:second-conjunct}
        \therefore \,
        \forall x \, \Big[ x \in A \implies \exists! y \, \big( y \in B \land \orderedpair{x}{y} \in f \big) \Big]
    \end{equation}
    Since \eqref{eq:first-conjunct} and \eqref{eq:second-conjunct} are provable from the assumption that \function{f}{A}{B},
    we conclude by the \aref{rule-of-inference:deduction-theorem}{Deduction Theorem} that,
    \begin{multline}
        \label{eq:f-A-B-implies-element-of-pairs-and-uniqueness-of-y}
        \function{f}{A}{B}
        \implies
        \bigg( \forall x \, \Big[ x \in f \implies \exists a,b \, \big( x = \orderedpair{a}{b} \land a \in A \land b \in B \big) \Big] \\
        \land
        \forall x \, \Big[ x \in A \implies \exists! y \, \big( y \in B \land \orderedpair{x}{y} \in f \big) \Big]
        \bigg)
    \end{multline}

    \vspace{0.2cm}
    \noindent
    To prove the converse, we start with the supposition that \( f \) is a set for which the following two statements hold true:
    \begin{gather}
        \forall x \, \Big[ x \in f \implies \exists a,b \, \big( x = \orderedpair{a}{b} \land a \in A \land b \in B \big) \Big]
        \label{eq:assumed-first-conjunct} \\
        \forall x \, \Big[ x \in A \implies \exists! y \, \big( y \in B \land \orderedpair{x}{y} \in f \big) \Big]
        \label{eq:assumed-second-conjunct}
    \end{gather}
    The goal is to prove that \eqref{eq:assumed-first-conjunct} and
    \eqref{eq:assumed-second-conjunct} do indeed imply \function{f}{A}{B}.
    This can be done by proving that \( f \) satisfies \cref{dfn:function-basic}
    and that \( \dom f = A \) and \( \ran f \subseteq B \).

    \vspace{0.3cm}
    \noindent
    Suppose that \( x_{0} \in f \). Thus, \( x_{0} = \orderedpair{\alpha}{\beta} \) for some \( \alpha \in A \) and \( \beta \in B \)
    by \eqref{eq:assumed-first-conjunct}. Clearly, \( x_{0} \) is an ordered pair, so,
    \begin{equation}
        \label{eq:first-conjunct-of-function-definition}
        \forall x \, \Big[ x \in f \implies \exists a,b \, \big( x = \orderedpair{a}{b} \big) \Big]
    \end{equation}
    Next, let \( x_{0} \), \( y_{0} \), and \( z_{0} \) be arbitrary objects such that \( \orderedpair{x_{0}}{y_{0}} \in f \)
    and \( \orderedpair{x_{0}}{z_{0}} \in f \). By \eqref{eq:assumed-first-conjunct}, this means that,
    \begin{equation*}
        \begin{aligned}
            \orderedpair{x_{0}}{y_{0}} = \orderedpair{\alpha}{\beta}& \\
            \orderedpair{x_{0}}{z_{0}} = \orderedpair{\gamma}{\delta}& \\
        \end{aligned}
        \begin{aligned}
            &\quad\text{for some \( \alpha \in A \) and \( \beta \in B \)}\\
            &\quad\text{for some \( \gamma \in A \) and \( \delta \in B \)}
        \end{aligned}
    \end{equation*}
    By \cref{thm:equality-of-ordered-pairs}, it follows that \( x_{0} = \alpha = \gamma \),
    \( y_{0} = \beta \), \( z_{0} = \delta \).
    By \aref{rule-of-inference:universal-instantiation}{Universal Instantiation} of \eqref{eq:assumed-second-conjunct}, we have,
    \begin{equation*}
        \alpha \in A \implies \exists! y \, \big( y \in B \land \orderedpair{\alpha}{y} \in f \big)
    \end{equation*}
    Since \( \alpha \in A \), we infer \( \exists! y \, \mleft( y \in B \land \orderedpair{\alpha}{y} \in f \mright) \) by
    \aref{rule-of-inference:modus-ponens}{Modus Ponens}. Equivalently,
    \begin{equation}
        \label{eq:re-express-uniqueness-equivalently}
        \exists y \, \Big[ \mleft( y \in B \land \orderedpair{\alpha}{y} \in f \mright) \land
        \forall z \, \big( \mleft( z \in B \land \orderedpair{\alpha}{z} \in f \mright) \implies y = z \big)\Big]
    \end{equation}
    Thus, by applying \aref{rule-of-inference:existential-instantiation}{Existential Instantiation} on \eqref{eq:re-express-uniqueness-equivalently}, we know that
    there exists some \( \mu \) such that \( \mu \in B \land \orderedpair{\alpha}{\mu} \in f \) and,
    \begin{equation}
        \label{eq:mu-is-unique}
        \forall z \, \Big[ \mleft( z \in B \land \orderedpair{\alpha}{z} \in f \mright) \implies \mu = z \Big]
    \end{equation}
    Since \( \orderedpair{\alpha}{y_{0}} \in f \) and \( \orderedpair{\alpha}{z_{0}} \in f \) where \( y_{0} \in B \)
    and \( z_{0} \in B \), we conclude by \aref{rule-of-inference:modus-ponens}{Modus Ponens} and \eqref{eq:mu-is-unique} that \( \mu = y_{0} = z_{0} \).
    We have just derived \( y_{0} = z_{0} \) from the supposition that
    \( \orderedpair{x_{0}}{y_{0}} \in f \land \orderedpair{x_{0}}{z_{0}} \in f \). Thus, by the \aref{rule-of-inference:deduction-theorem}{Deduction Theorem}, we
    conclude,
    \begin{equation*}
        \orderedpair{x_{0}}{y_{0}} \in f \land \orderedpair{x_{0}}{z_{0}} \in f \implies y_{0} = z_{0}
    \end{equation*}
    Since \( x_{0} \), \( y_{0} \), and \( z_{0} \) are arbitrary, \aref{rule-of-inference:universal-generalization}{Universal Generalization} yields,
    \begin{equation}
        \label{eq:second-conjunct-of-function-definition}
        \forall x \, \forall y \, \forall z \, \Big[ \orderedpair{x}{y} \in f \land \orderedpair{x}{z} \in f \implies y = z \Big]
    \end{equation}
    With both \eqref{eq:first-conjunct-of-function-definition} and \eqref{eq:second-conjunct-of-function-definition} being true,
    it follows that \( f \) is indeed a function (see \cref{dfn:function-basic}).

    \vspace{0.3cm}
    \noindent
    Proving that \( \dom f = A \) and \( \ran f \subseteq B \)
    is quite straightforward. Suppose, for instance, that \( x_{0} \in A \).
    Then by \eqref{eq:assumed-second-conjunct}, it follows that
    \( \orderedpair{x_{0}}{y_{0}} \in f \) for some unique \( y_{0} \in B \).
    Thus, \( x_{0} \in \dom f \) by \cref{dfn:domain-of-a-set-of-ordered-pairs}.
    Conversely, if \( x_{0} \in \dom f \), then \( \orderedpair{x_{0}}{y_{0}} \)
    for some \( y_{0} \). Since \( \orderedpair{x_{0}}{y_{0}} \in f \), it follows
    from \eqref{eq:assumed-first-conjunct} that
    \( \orderedpair{x_{0}}{y_{0}} = \orderedpair{\alpha}{\beta} \) for some
    \( \alpha \in A \) and \( \beta \in B \). By \cref{thm:equality-of-ordered-pairs},
    this means that \( x_{0} = \alpha \) and \( y_{0} = \beta \). Since \( \alpha \in A \)
    and \( x_{0} \in \alpha \), necessarily \( x_{0} \in A \). In other words,
    \begin{equation*}
        \begin{aligned}
            x_{0} \in \dom f \iff x_{0} \in A&
        \end{aligned}
        \begin{aligned}
            &\text{for any \( x_{0} \)}
        \end{aligned}
    \end{equation*}
    By \aref{rule-of-inference:universal-generalization}{Universal Generalization} and \nameref{axm:axiom-of-extensionality}, it follows that,
    \begin{equation*}
        \dom f = A
    \end{equation*}
    We can use a similar argument to show that \( \ran f \subseteq B \). Let \( y_{0} \)
    be an arbitrary object such that \( y_{0} \in \ran f \). Thus \( \orderedpair{x_{0}}{y_{0}} \in f \)
    for some \( x_{0} \) by \cref{dfn:range-of-a-set-of-ordered-pairs}. But if \( \orderedpair{x_{0}}{y_{0}} \in f \),
    then by \eqref{eq:assumed-first-conjunct}, there must exist some \( \alpha \in A \) and \( \beta \in B \) such
    that \( \orderedpair{x_{0}}{y_{0}} = \orderedpair{\alpha}{\beta} \), which in turn implies \( x_{0} \in A \)
    and \( y_{0} \in B \). Thus,
    \begin{equation*}
        \begin{aligned}
            y_{0} \in \ran f \implies y_{0} \in B&
        \end{aligned}
        \begin{aligned}
            &\text{for any \( y_{0} \)}
        \end{aligned}
    \end{equation*}
    By \aref{rule-of-inference:universal-generalization}{Universal Generalization} and \cref{dfn:subset}, necessarily \( \ran f \subseteq B \). Since \( f \) is a
    function with \( \dom f = A \) and \( \ran f \subseteq B \), it follows that,
    \begin{equation*}
        \function{f}{A}{B}
    \end{equation*}
    Thus, having derived \function{f}{A}{B} from \eqref{eq:assumed-first-conjunct} and \eqref{eq:assumed-second-conjunct},
    we conclude by the \aref{rule-of-inference:deduction-theorem}{Deduction Theorem},
    \begin{multline}
        \label{eq:element-of-pairs-and-uniqueness-of-y-implies-f-A-B}
        \bigg(
        \forall x \, \Big[ x \in f \implies \exists a,b \, \big( x = \orderedpair{a}{b} \land a \in A \land b \in B \big) \Big] \\
        \land
        \forall x \, \Big[ x \in A \implies \exists! y \, \big( y \in B \land \orderedpair{x}{y} \in f \big) \Big]
        \bigg)
        \implies
        \function{f}{A}{B}
    \end{multline}
    Together, \eqref{eq:f-A-B-implies-element-of-pairs-and-uniqueness-of-y} and \eqref{eq:element-of-pairs-and-uniqueness-of-y-implies-f-A-B}
    establish the required biconditional.
\end{proof}

\begin{thm}
    \label{thm:f-is-a-function-with-empty-domain-and-empty-range-if-and-only-if-it-is-the-empty-set}
    For any set \( f \), \( \function{f}{\emptyset}{\emptyset} \) if and only if \( f = \emptyset \).
    \begin{equation*}
        \forall f \, \Big[ \function{f}{\emptyset}{\emptyset} \iff f = \emptyset \Big]
    \end{equation*}
\end{thm}

\begin{proof}
    Let \( f \) be a set such that \( \function{f}{\emptyset}{\emptyset} \). By \cref{thm:full-definition-of-the-arrow-notation},
    \begin{equation}
        \label{eq:any-element-of-f-must-be-an-ordered-pair-from-empty-and-empty}
        \forall x \, \Big[ x \in f \implies \exists a,b \, \big( x = \orderedpair{a}{b} \land a \in \emptyset \land b \in \emptyset \big) \Big]
    \end{equation}
    Let \( f \) be an arbitrary object.
    First, suppose that \( f \neq \emptyset \). This means that there exists an \( x_{0} \) such that \( x_{0} \in f \).
    But by \eqref{eq:any-element-of-f-must-be-an-ordered-pair-from-empty-and-empty}, this would mean that \( x_{0} = \orderedpair{\alpha}{\beta} \)
    for some \( \alpha \in \emptyset \) and \( \beta \in \emptyset \), which is an obvious contradiction because \( \emptyset \) cannot contain
    any elements by definition. Thus, the supposition that \( f \neq \emptyset \) must be false by \aref{rule-of-inference:reductio-ad-absurdum}{Reductio ad Absurdum}.
    \begin{equation}
        \label{eq:f-is-a-function-with-empty-domain-and-empty-range-implies-f-equals-empty}
        \therefore \, \function{f}{\emptyset}{\emptyset} \implies f = \emptyset
    \end{equation}
    Conversely, suppose that \( f = \emptyset \). But if \( f = \emptyset \), then both conjuncts in \cref{thm:full-definition-of-the-arrow-notation}
    are vacuously true so that \( \function{f}{\emptyset}{\emptyset} \) is also true.
    \begin{equation}
        \label{eq:f-equals-empty-implies-f-is-a-function-with-empty-domain-and-empty-range}
        f = \emptyset \implies \function{f}{\emptyset}{\emptyset}
    \end{equation}
    Having derived both \eqref{eq:f-is-a-function-with-empty-domain-and-empty-range-implies-f-equals-empty} and
    \eqref{eq:f-equals-empty-implies-f-is-a-function-with-empty-domain-and-empty-range}, we can conclude,
    \begin{equation*}
        \function{f}{\emptyset}{\emptyset} \iff f = \emptyset \qedhere
    \end{equation*}
\end{proof}

\begin{thm}[Equality]
    \label{thm:equality-of-functions}
    If \( f \) and \( g \) are functions, then \( f = g \) if and only if,
    \begin{equation*}
        \Big( \dom f = \dom g \Big) \land \forall x \, \Big( x \in \dom f \implies f(x) = g(x) \Big)
    \end{equation*}
\end{thm}

\begin{proof}
    Let \( f \) and \( g \) be arbitrary functions. If \( f = g \), then by the \nameref{axm:axiom-of-extensionality},
    \begin{equation}
        \label{eq:f-and-g-have-the-same-elements}
        \forall x \, \mleft( x \in f \iff x \in g \mright)
    \end{equation}
    Let \( x_{0} \in \dom f \) be arbitrary. By \cref{dfn:domain-of-a-set-of-ordered-pairs}, clearly
    \( \orderedpair{x_{0}}{y_{0}} \in f \) for some \( y_{0} \). Thus, \( \orderedpair{x_{0}}{y_{0}} \in g \)
    by \eqref{eq:f-and-g-have-the-same-elements} so that \( \exists y \, \big( \orderedpair{x_{0}}{y} \in g \big) \).
    By \cref{dfn:domain-of-a-set-of-ordered-pairs}, it follows that \( x_{0} \in \dom g \).
    Note that \( f(x_{0}) = y_{0} = g(x_{0}) \).
    \begin{equation}
        \label{eq:forward-implication-first-half}
        \therefore \, \forall x \, \Big( x \in \dom f \implies x \in \dom g \Big)
        \land
        \forall x \, \Big( x \in \dom f \implies f(x) = g(x) \Big)
    \end{equation}
    On the other hand, let \( x_{0} \in \dom g \) be arbitrary. Once again, \( \orderedpair{x_{0}}{y_{0}} \in g \)
    for some \( y_{0} \). And by \eqref{eq:f-and-g-have-the-same-elements}, necessarily \( \orderedpair{x_{0}}{y_{0}} \in f \)
    so that \( \exists y \, \big( \orderedpair{x_{0}}{y} \in f \big) \).
    \begin{equation}
        \label{eq:forward-implication-second-half}
        \therefore \, \forall x \, \Big( x \in \dom g \implies x \in \dom f \Big)
    \end{equation}
    Combining \eqref{eq:forward-implication-first-half} and \eqref{eq:forward-implication-second-half} yields,
    \begin{equation*}
        \Big( \dom f = \dom g \Big) \land \forall x \, \Big( x \in \dom f \implies f(x) = g(x) \Big)
    \end{equation*}
    Since the statement was derived on the assumption that \( f = g \), we conclude by the \aref{rule-of-inference:deduction-theorem}{Deduction Theorem} that,
    \begin{equation}
        \label{eq:f-equals-g-implies-equal-domains-and-equal-ranges}
        \Big( f = g \Big) \implies \Big( \big( \dom f = \dom g \big) \land \forall x \, \big( x \in \dom f \implies f(x) = g(x) \big) \Big)
    \end{equation}
    To prove the converse, suppose that,
    \begin{gather}
        \dom f = \dom g \label{eq:suppose-equal-domain} \\
        \forall x \, \big( x \in \dom f \implies f(x) = g(x) \big) \label{eq:suppose-equal-values}
    \end{gather}
    Let \( x_{0} \in f\) be arbitrary.
    Since \( f \) is a function, it follows from \cref{dfn:function-basic} that
    \( x_{0} = \orderedpair{a_{0}}{b_{0}} \) for some \( a_{0} \) and \( b_{0} \).
    Applying \aref{rule-of-inference:universal-instantiation}{Universal Instantiation} on \eqref{eq:suppose-equal-values} yields,
    \begin{equation*}
        a_{0} \in \dom f \implies f(a_{0}) = g(a_{0})
    \end{equation*}
    Clearly \( a_{0} \in \dom f\), so \( f(a_{0}) = g(a_{0}) \) by \aref{rule-of-inference:modus-ponens}{Modus Ponens}.
    Since \( \orderedpair{a_{0}}{b_{0}} \in f \), it follows from
    \cref{dfn:functional-notation} that \( b_{0} = f(a_{0}) \). Thus, \( g(a_{0}) = f(a_{0}) = b_{0} \).
    Since \( b_{0} = g(a_{0}) \), it follows from \cref{dfn:functional-notation} that \( x_{0} = \orderedpair{a_{0}}{b_{0}} \in g \).
    \begin{equation}
        \label{eq:f-is-subset-of-g}
        \therefore \, x_{0} \in f \implies x_{0} \in g
    \end{equation}
    On the other hand, suppose that \( x_{0} \in g \). Again, \( g \) is a function, so \( x_{0} = \orderedpair{a_{0}}{b_{0}}\) for some
    \( a_{0} \) and \( b_{0} \) with \( a_{0} \in \dom g \) and \( b_{0} = g(a_{0}) \). But \( \dom f = \dom g \) by \eqref{eq:suppose-equal-domain},
    so \( a_{0} \in \dom g \) implies \( a_{0} \in \dom f \) by the \nameref{axm:axiom-of-extensionality}. Since \( a_{0} \in \dom f \), we have
    \( f(a_{0}) = g(a_{0}) = b_{0} \) by \aref{rule-of-inference:modus-ponens}{Modus Ponens} and \eqref{eq:suppose-equal-values}. In other words, \( x_{0} = \orderedpair{a_{0}}{b_{0}} \in f \).
    \begin{equation}
        \label{eq:g-is-subset-of-f}
        \therefore \, x_{0} \in f \implies x_{0} \in f
    \end{equation}
    From \eqref{eq:f-is-subset-of-g} and \eqref{eq:g-is-subset-of-f}, we infer the biconditional,
    \begin{equation*}
        x_{0} \in f \iff x_{0} \in g
    \end{equation*}
    Since \( x_{0} \) is arbitrary, we can generalize,
    \begin{equation*}
        \forall x \, \big( x \in f \iff x \in g \big)
    \end{equation*}
    And by the \nameref{axm:axiom-of-extensionality}, \( f = g \). Given that \( f = g \) is provable from the assumption that
    \( \dom f = \dom g \) and \( \forall x \, \big( x \in \dom f \implies f(x) = g(x) \big) \), we conclude by the Deduction
    Theorem that,
    \begin{equation}
        \label{eq:equal-domains-and-equal-ranges-implies-f-equals-g}
        \Big( \dom f = \dom g \Big) \land \forall x \, \Big( x \in \dom f \implies f(x) = g(x) \Big) \implies \Big( f = g \Big)
    \end{equation}
    Together, \eqref{eq:f-equals-g-implies-equal-domains-and-equal-ranges} and \eqref{eq:equal-domains-and-equal-ranges-implies-f-equals-g}
    establish the required biconditional.
\end{proof}

\begin{thm}[Composite Functions]
    \label{thm:the-composition-of-any-two-functions-is-also-a-function}
    If \( f \) and \( g \) are functions, then \composition{g}{f} is also a function. Furthermore:
    \begin{enumerate}[topsep=2pt,noitemsep,label=(\roman*)]
        \item \( \dom \mleft( \composition{g}{f} \mright) = \imageSet{f^{-1}}{\dom g} \) \label{enum:composite-function-claim-1}
        \item \( \forall x \, \big( x \in \dom \mleft( \composition{g}{f} \mright) \implies \mleft( \composition{g}{f} \mright)(x) = g(f(x)) \big) \)
        \label{enum:composite-function-claim-2}
    \end{enumerate}
\end{thm}

\begin{proof}
    Let \( f \) and \( g \) be arbitrary functions.
    By \cref{dfn:composition-of-two-sets-of-ordered-pairs},
    \composition{g}{f} is clearly a set of ordered pairs.
    To prove that \composition{g}{f} is a function, let \( x_{0} \) be an arbitrary object such that
    \( \orderedpair{x_{0}}{y_{1}} \in \composition{g}{f} \) and \( \orderedpair{x_{0}}{y_{2}} \in \composition{g}{f} \)
    for some \( y_{1} \) and \( y_{2} \). Thus, there exists some \( z_{1} \) and \( z_{2} \) such that,
    \begin{gather*}
        \orderedpair{x_{0}}{z_{1}} \in f \land \orderedpair{z_{1}}{y_{1}} \in g \\
        \orderedpair{x_{0}}{z_{2}} \in f \land \orderedpair{z_{2}}{y_{2}} \in g
    \end{gather*}
    Since \( f \) is a function, \( \orderedpair{x_{0}}{z_{1}} \in f \land \orderedpair{x_{0}}{z_{2}} \in f \) implies
    \( z_{1} = z_{2} \). Since \( g \) is also a function, \( z_{1} = z_{2} \) then implies \( y_{1} = y_{2} \). In other words,
    \begin{equation*}
        \orderedpair{x_{0}}{y_{1}} \in \composition{g}{f} \land \orderedpair{x_{0}}{y_{2}} \in \composition{g}{f}
        \implies
        y_{1} = y_{2}
    \end{equation*}
    Since \( x_{0} \), \( y_{1} \), and \( y_{2} \) are arbitrary, we have by \aref{rule-of-inference:universal-generalization}{Universal Generalization},
    \begin{equation*}
        \forall x \, \forall y \, \forall z \,
        \Big( \orderedpair{x}{y} \in \composition{g}{f} \land \orderedpair{x}{z} \in \composition{g}{f}
        \implies y = z \Big)
    \end{equation*}
    By \cref{dfn:function-basic}, \composition{g}{f} is indeed a function.

    \vspace{0.3cm}
    \noindent
    To prove \ref{enum:composite-function-claim-1}, let \( x_{0} \) be an arbitrary object.
    It is easy to show that \( x_{0} \in \dom \mleft( \composition{g}{f} \mright) \) and
    \( x_{0} \in \imageSet{f}{\dom g} \) are indeed equivalent statements.
    We proceed by applying
    \cref{dfn:domain-of-a-set-of-ordered-pairs,dfn:inverse-of-a-set-of-ordered-pairs,dfn:image-of-a-set-of-ordered-pairs}:
    \begin{align*}
        &x_{0} \in \dom \mleft( \composition{g}{f} \mright) \\
        &\iff
        \exists y \, \Big[ \orderedpair{x_{0}}{y} \in \composition{g}{f} \Big] \\
        &\iff
        \exists y \, \Big[ \exists z \, \big( \orderedpair{x_{0}}{z} \in f \land \orderedpair{z}{y} \in g \big) \Big] \\
        &\iff
        \exists z \, \Big[ \orderedpair{x_{0}}{z} \in f \land \exists y \, \big( \orderedpair{z}{y} \in g \big) \Big] \\
        &\iff
        \exists z \, \Big[ \orderedpair{x_{0}}{z} \in f \land z \in \dom g \Big] \\
        &\iff
        \exists z \, \Big[ z \in \dom g \land \orderedpair{x_{0}}{z} \in f \Big] \\
        &\iff
        \exists z \, \Big[ z \in \dom g \land \orderedpair{z}{x_{0}} \in f^{-1} \Big] \\
        &\iff
        x_{0} \in \imageSet{f^{-1}}{\dom g}
    \end{align*}
    By \aref{rule-of-inference:universal-generalization}{Universal Generalization} and the \nameref{axm:axiom-of-extensionality},
    we conclude,
    \begin{equation*}
        \dom \mleft( \composition{g}{f} \mright) = \imageSet{f^{-1}}{\dom g}
    \end{equation*}

    \vspace{0.3cm}
    \noindent
    To prove \ref{enum:composite-function-claim-2}, we let \( x_{0} \in \dom \mleft( \composition{g}{f} \mright) \) be arbitrary.
    By \cref{dfn:domain-of-a-set-of-ordered-pairs}, it follows that \( \orderedpair{x_{0}}{y_{0}} \in \composition{g}{f} \) for
    some \( y_{0} \). Also, \( \orderedpair{x_{0}}{z_{0}} \in f \) and \( \orderedpair{z_{0}}{y_{0}} \in g \) for some \( z_{0} \).
    Applying \cref{dfn:functional-notation}, it is clear that \( z_{0} = f(x_{0}) \) and
    \( y_{0} = g(z_{0}) = g(f(x_{0})) = \mleft( \composition{g}{f} \mright)(x_{0}) \).
\end{proof}

\begin{thm}
    \label{thm:a-function-is-invertible-if-and-only-if-it-is-injective}
    A function \( f \) is invertible if and only if it is injective.
\end{thm}

\begin{proof}
    Let \( f \) be an arbitrary function. Suppose that \( f \) is injective. That is,
    \begin{equation}
        \label{eq:f-is-injective}
        \forall x \, \forall y \, \big[ x \in \dom f \land y \in \dom f \implies \big( f(x) = f(y) \implies x = y \big) \big]
    \end{equation}
    Clearly, \( f^{-1} \) exists and is a set of ordered pairs. Let \( x_{0} \), \( y_{0} \), and \( z_{0} \) be arbitrary
    objects such that \( \orderedpair{x_{0}}{y_{0}} \in f^{-1} \) and \( \orderedpair{x_{0}}{z_{0}} \in f^{-1} \). By,
    \cref{dfn:inverse-of-a-set-of-ordered-pairs}, this implies \( \orderedpair{y_{0}}{x_{0}} \in f \) and \( \orderedpair{z_{0}}{x_{0}} \in f \).
    This also means that \( y_{0} \in \dom f \) and \( z_{0} \in \dom f \). But \( f \) is injective, so we have by \eqref{eq:f-is-injective},
    \begin{equation*}
        f(y_{0}) = f(z_{0}) \implies y_{0} = z_{0}
    \end{equation*}
    But \( f(y_{0}) = x_{0} = f(z_{0}) \), so by \aref{rule-of-inference:modus-ponens}{Modus Ponens}, we infer \( y_{0} = z_{0} \).
    \begin{equation*}
        \therefore \, \forall x \, \forall y \, \forall z \, \big[ \orderedpair{x}{y} \in f^{-1} \land
        \orderedpair{x}{y} \in f^{-1} \implies y = z \big]
    \end{equation*}
    Thus, \( f^{-1} \) is a function, so \( f \) is invertible by \cref{dfn:invertibility-of-functions}.

    \vspace{0.3cm}
    \noindent
    To prove the converse, suppose that  \( f \) is invertible (that is, \( f^{-1} \) is a function).
    Let \( x_{1} \) and \( x_{2} \) be arbitrary objects such that \( x_{1},x_{2} \in \dom f \)
    and \( f(x_{1}) = f(x_{2}) \). By
    \cref{dfn:functional-notation}, it follows that,
    \begin{gather*}
        \orderedpair{x_{1}}{f(x_{1})} \in f \\
        \orderedpair{x_{2}}{f(x_{2})} \in f
    \end{gather*}
    And by \cref{dfn:inverse-of-a-set-of-ordered-pairs}, it follows that,
    \begin{gather*}
        \orderedpair{f(x_{1})}{x_{1}} \in f^{-1} \\
        \orderedpair{f(x_{2})}{x_{2}} \in f^{-1}
    \end{gather*}
    Since \( f^{-1} \) is a function and \( f(x_{1}) = f(x_{2}) \), necessarily \( x_{1} = x_{2} \).
    Thus, we have by the \aref{rule-of-inference:deduction-theorem}{Deduction Theorem} \( (f(x_{1}) = f(x_{2})) {\implies} (x_{1} = x_{2}) \)
    for any \( x_{1}, x_{2} \in \dom f \).
    Finally, \aref{rule-of-inference:universal-generalization}{Universal Generalization} yields,
    \begin{equation*}
        \forall x \, \forall y \, \Big[ x \in \dom f \land y \in \dom f \implies \Big( f(x) = f(y) \implies x = y \Big) \Big]
    \end{equation*}
    Thus, \( f \) is injective by \cref{dfn:injection}.
\end{proof}

\begin{cor}
    If a function \( f \) is invertible, then \( f^{-1} \) is also invertible because \( \mleft( f^{-1} \mright)^{-1} = f \) by
    \cref{thm:inverse-of-the-inverse-of-a-set-of-ordered-pairs-is-the-set-itself}.
\end{cor}

\begin{thm}
    \phantomsection \label{thm:the-inverse-of-a-bijection-is-also-a-bijection}
    For any two sets \( A \) and \( B \), if \( \bijection{f}{A}{B} \), then \( \bijection{f^{-1}}{B}{A} \).
\end{thm}

\begin{proof}
    Let \( A \) and \( B \) be two arbitrary sets and \( f \) be a bijection between them.
    By \cref{dfn:inverse-of-a-set-of-ordered-pairs}, it follows that,
    \begin{equation}
        \phantomsection \label{eq:f-inverse-in-setbuilder-form}
        f^{-1} = \setbuilder{\orderedpair{x}{y}}{\orderedpair{y}{x} \in f}
    \end{equation}
    Since \( \dom f = A \) and \( \ran f = B \), we have by \cref{thm:domain-is-the-range-of-inverse-and-vice-versa}
    \( \dom f^{-1} = B \) and \( \ran f^{-1} = A \). Furthermore, \( f \) is a \aref{dfn:bijection}{bijection}, so it must also be \aref{dfn:injection}{injective}.
    Thus, \cref{thm:a-function-is-invertible-if-and-only-if-it-is-injective} guarantees that \( f^{-1} \) is a function.
    In other words, \( \surjection{f^{-1}}{B}{A} \).

    \vspace{0.3cm}
    \noindent
    To show that \( f^{-1} \) is also injective, we let \( x_{1}, x_{2} \in B \) and \( y \in A \) be arbitrary objects such that
    \( \orderedpair{x_{1}}{y} \in f^{-1} \) and \( \orderedpair{x_{2}}{y} \in f^{-1} \). By \eqref{eq:f-inverse-in-setbuilder-form},
    it follows that \( \orderedpair{y}{x_{1}} \in f \) and \( \orderedpair{y}{x_{2}} \in f \). But \( f \) is a function,
    so \cref{dfn:function-basic} guarantees that \( x_{1} = x_{2} \).
    Thus, \( \orderedpair{x_{1}}{y} \in f^{-1} \,\,{\land}\,\, \orderedpair{x_{2}}{y} \in f^{-1} {\implies} x_{1} = x_{2} \) for any
    arbitrary \( x_{1},x_{2} \in B \) and \( y \in A \).
    By \cref{dfn:injection}, it follows that \( f^{-1} \) is injective.
\end{proof}

\begin{thm}
    \phantomsection \label{thm:for-any-function-f-the-image-of-A-under-f-inverse-is-a-subset-of-A}
    For any function \( f \) and any arbitrary set \( A \), \( \imageSet{f}{\imageSet{f^{-1}}{A}} \subseteq A \).
\end{thm}

\begin{proof}
    Let \( f \) be a function and \( A \) be an arbitrary set. Let \( S \) denote the set \( \imageSet{f^{-1}}{A} \).
    Thus, we have by \cref{dfn:image-of-a-set-of-ordered-pairs,dfn:inverse-of-a-set-of-ordered-pairs},
    \begin{equation}
        \label{eq:S-is-the-image-of-A-under-f-inverse}
        S = \imageSet{f^{-1}}{A} = \setbuilderbig{y}{\exists x \, \mleft( x \in A \,\,{\land}\,\, \orderedpair{y}{x} \in f \mright)}
    \end{equation}
    We now proceed to demonstrate that \( \imageSet{f}{S} \subseteq A \).

    \vspace{0.3cm}
    \noindent
    Let \( z_{0} \in \imageSet{f}{S} \) be arbitrary. By \cref{dfn:image-of-a-set-of-ordered-pairs}, it follows that \( \orderedpair{x_{0}}{z_{0}} \in f \)
    for some \( x_{0} \in S \).
    Since \( x_{0} \in S \), we can then infer from \eqref{eq:S-is-the-image-of-A-under-f-inverse} that \( \orderedpair{x_{0}}{y_{0}} \in f \) for some
    \( y_{0} \in A \). In other words,
    \begin{equation*}
        \orderedpair{x_{0}}{z_{0}} \in f \,\,{\land}\,\, \orderedpair{x_{0}}{y_{0}} \in f
    \end{equation*}
    But \( f \) is a function, so it follows from \cref{dfn:function-basic} that \( z_{0} = y_{0} \). Since \( y_{0} \in A \) and \( y_{0} = z_{0} \),
    we have by \aref{rule-of-inference:substitution}{Substitution} that \( z_{0} \in A \). By the \aref{rule-of-inference:deduction-theorem}{Deduction Theorem},
    we conclude that \( z_{0} \in \imageSet{f}{S} {\implies} z_{0} \in A \) for any arbitrary \( z_{0} \).
    Finally, we conclude by \aref{rule-of-inference:universal-generalization}{Universal Generalization} that \( \imageSet{f}{S} \subseteq A \).
\end{proof}

\begin{thm}
    \phantomsection \label{thm:image-of-f-inverse-A-under-f-is-A-if-f-is-a-bijection-and-A-subset-of-range-f}
    If \( f \) is a bijection and \( A \subseteq \ran f \), then \( \imageSet{f}{\imageSet{f^{-1}}{A}} = A \).
\end{thm}

\begin{proof}
    Let \( f \) be a \aref{dfn:bijection}{bijective function} and \( A \) be any set such that \( A \subseteq \ran f \).
    Since \( f \) is a function, it follows from \cref{thm:for-any-function-f-the-image-of-A-under-f-inverse-is-a-subset-of-A} that
    \( \imageSet{f}{\imageSet{f^{-1}}{A}} \subseteq A \). Thus, it remains to show that \( A \subseteq \imageSet{f}{\imageSet{f^{-1}}{A}} \).
    We start by writing down the definition of \( \imageSet{f}{\imageSet{f^{-1}}{A}} \)
    according to \cref{dfn:image-of-a-set-of-ordered-pairs,dfn:inverse-of-a-set-of-ordered-pairs}:
    \begin{equation}
        \label{eq:image-of-f-under-f-inverse-A-in-setbuilder-notation}
        \imageSet{f}{\imageSet{f^{-1}}{A}} = \setbuilderBig{y}{\exists x \, \exists z \, \big[ z \in A \,\,{\land}\,\, \orderedpair{x}{z} \in f
        \,\,{\land}\,\, \orderedpair{x}{y} \in f \big]}
    \end{equation}
    Now let \( y \in A \) be arbitrary. Since \( A \subseteq \ran f \) by \aref{rule-of-inference:assumption}{Assumption}, it follows that
    \( y \in \ran f \) and thus \( \orderedpair{x}{y} \in f \) for some \( x \in \dom f \) by \cref{dfn:range-of-a-set-of-ordered-pairs}.
    But \( \orderedpair{x}{y} \in f {\iff} \orderedpair{x}{y} \in f \,\,{\land}\,\, \orderedpair{x}{y} \in f \) is a tautology, and since
    \( y \in A \), we have by \aref{rule-of-inference:conjunction-introduction}{Conjunction Introduction}:
    \begin{equation*}
        y \in A \,\,{\land}\,\, \orderedpair{x}{y} \in f \,\,{\land}\,\, \orderedpair{x}{y} \in f
    \end{equation*}
    In other words, the formula \( \exists x \, \exists z \, \mleft[ z \in A \,\,{\land}\,\, \orderedpair{x}{z} \in f \,\,{\land}\,\, \orderedpair{x}{y} \in f \mright] \)
    is satisfied and thus \( y \in \imageSet{f}{\imageSet{f^{-1}}{A}} \) by \eqref{eq:image-of-f-under-f-inverse-A-in-setbuilder-notation}.
    By the \aref{rule-of-inference:deduction-theorem}{Deduction Theorem}, we have for any arbitrary \( y \),
    \begin{equation*}
        y \in A \implies y \in \imageSet{f}{\imageSet{f^{-1}}{A}}
    \end{equation*}
    Thus, \( A \subseteq \imageSet{f}{\imageSet{f^{-1}}{A}} \) follows readily from
    \aref{rule-of-inference:universal-generalization}{Universal Generalization} and \cref{dfn:subset}.
    Since \( \imageSet{f}{\imageSet{f^{-1}}{A}} \subseteq A \) and \( A \subseteq \imageSet{f}{\imageSet{f^{-1}}{A}} \) are both true,
    necessarily \( \imageSet{f}{\imageSet{f^{-1}}{A}} = A \).
\end{proof}

\begin{thm}
    \label{thm:the-composition-of-injective-functions-is-injective}
    If \( f \) and \( g \) are injective functions, then \composition{g}{f} is also injective.
\end{thm}

\begin{proof}
    Let \( f \) and \( g \) be injective functions. It follows from \cref{thm:the-composition-of-any-two-functions-is-also-a-function} that
    \composition{g}{f} is a function. Let \( x_{1} \) and \( x_{2} \) be elements of \( \dom \mleft( \composition{g}{f} \mright) \) such that
    \( \mleft( \composition{g}{f} \mright)(x_{1}) = \mleft( \composition{g}{f} \mright)(x_{2}) = y \). By \cref{dfn:functional-notation},
    we have,
    \begin{gather*}
        \orderedpair{x_{1}}{y} \in \composition{g}{f} \\
        \orderedpair{x_{2}}{y} \in \composition{g}{f}
    \end{gather*}
    Applying \cref{dfn:composition-of-two-sets-of-ordered-pairs}, we infer,
    \begin{equation*}
        \begin{aligned}
            \orderedpair{x_{1}}{z_{1}} \in f \land \orderedpair{z_{1}}{y} \in g&\\
            \orderedpair{x_{2}}{z_{2}} \in f \land \orderedpair{z_{2}}{y} \in g&
        \end{aligned}
        \text{  }
        \begin{aligned}
            &\quad\text{for some \( z_{1} \)}\\
            &\quad\text{for some \( z_{2} \)}
        \end{aligned}
    \end{equation*}
    Clearly, \( g(z_{1}) = g(z_{2}) = y \). Since \( g \) is injective, necessarily \( z_{1} = z_{2} \). Since \( f \) is also injective
    and \( f(x_{1}) = z_{1} = z_{2} = f(x_{2}) \), it follows that \( x_{1} = x_{2} \). Thus,
    \begin{equation*}
        x_{1} \in \dom \mleft( \composition{g}{f} \mright)
        \land
        x_{2} \in \dom \mleft( \composition{g}{f} \mright)
        \implies
        \Big( \mleft( \composition{g}{f} \mright)(x_{1}) = \mleft( \composition{g}{f} \mright)(x_{2}) \implies x_{1} = x_{2} \Big)
    \end{equation*}
    Since \( x_{1} \) and \( x_{2} \) are arbitrary, we conclude by \aref{rule-of-inference:universal-generalization}{Universal Generalization},
    \begin{equation*}
        \forall x \, \forall y \, \Big[ x \in \mleft( \composition{g}{f} \mright) \land y \in \mleft( \composition{g}{f} \mright)
        \implies
        \Big( \mleft( \composition{g}{f} \mright)(x) = \mleft( \composition{g}{f} \mright)(y) \implies x = y \Big) \Big]
    \end{equation*}
    Thus, \composition{g}{f} is injective by \cref{dfn:injection}.
\end{proof}

\begin{thm}
    \label{thm:for-injective-functions-the-image-of-intersection-is-equal-to-the-intersection-of-images}
    If \( f \) is an injective function and \( S \) is a non-empty system of sets, then,
    \begin{equation*}
        \imageSet{f}{\Intersect S} = \Intersect \setbuilderBig{t}{\exists z \, \big( z \in S \land t = \imageSet{f}{z} \big)}
    \end{equation*}
\end{thm}

\begin{proof}
    Let \( f \) be an injective function and \( S\) be a non-empty system of sets.
    Let the set \( T \) be defined as,
    \begin{equation*}
        T = \Intersect \setbuilderBig{t}{\exists z \, \big( z \in S \land t = \imageSet{f}{z} \big)}
    \end{equation*}
    It follows from \cref{thm:image-of-an-intersection-is-a-subset-of-the-intersection-of-all-individual-images} that
    \begin{equation*}
        \imageSet{f}{\Intersect S} \subseteq \Intersect T
    \end{equation*}
    Thus, it remains to prove the inclusion in the opposite direction.
    Let \( y_{0} \) be an arbitrary object such that \( y_{0} \in \Intersect T \).
    \begin{equation}
        \label{eq:y0-belongs-to-all-members-of-T}
        \therefore \, \forall t \, \big[ t \in T \implies y_{0} \in t \big]
    \end{equation}
    By the laws of logic, we can re-express \eqref{eq:y0-belongs-to-all-members-of-T} in other equivalent forms,
    \begin{align}
        \label{eq:y0-belongs-to-all-memebers-of-T-final-form}
        &\forall t \, \big[ t \in T \implies y_{0} \in t \big] \notag \\
        &\iff
        \forall t \, \big[ \exists z \, \big( z \in S \land t = \imageSet{f}{z} \big) \implies y_{0} \in t \big] \notag \\
        &\iff
        \forall t \, \big[ \lnot \, \exists z \, \big( z \in S \land t = \imageSet{f}{z} \big) \lor  y_{0} \in t \big] \notag \\
        &\iff
        \forall t \, \forall z \, \big[ z \notin S \lor t \neq \imageSet{f}{z} \lor y_{0} \in t \big] \notag \\
        &\iff
        \forall z \, \big[ z \notin S \lor \forall t \, \big( t = \imageSet{f}{z} \implies y_{0} \in t \big) \big] \notag \\
        &\iff
        \forall z \, \big[ z \notin S \lor y_{0} \in \imageSet{f}{z} \big] \notag \\
        &\iff
        \forall z \, \mleft[ z \in S \implies y_{0} \in \imageSet{f}{z} \mright]
    \end{align}
    Since \( S \neq \emptyset \), there exists an \( r \in S \). And by \eqref{eq:y0-belongs-to-all-memebers-of-T-final-form},
    necessarily \( y_{0} \in \imageSet{f}{r} \). In other words,
    \begin{equation}
        \label{eq:fxr-equals-y0}
        \begin{aligned}
            \orderedpair{x_{r}}{y_{0}} \in f &
        \end{aligned}
        \begin{aligned}
            &\text{for some \( x_{r} \in r \)}
        \end{aligned}
    \end{equation}
    Thus, if we can prove that \( x_{r} \in \Intersect S\), then we will have proven that \( y_{0} \in \imageSet{f}{\Intersect S} \).

    \vspace{0.5cm}
    \noindent
    To prove that \( x_{r} \in \Intersect S \), we let \( z_{0} \) be an arbitrary object. If \( z_{0} \in S \),
    then by \eqref{eq:y0-belongs-to-all-memebers-of-T-final-form}, we have \( y_{0} \in \imageSet{f}{z_{0}} \).
    That means, \( \orderedpair{x_{0}}{y_{0}} \in f \) for some \( x_{0} \in z_{0} \). By \cref{dfn:functional-notation},
    this also means that \( f(x_{0}) = y_{0} \). But by \eqref{eq:fxr-equals-y0}, \( f(x_{r}) \) is also equal to \( y_{0} \).
    And since \( f \) is injective, we have,
    \begin{equation*}
        f(x_{0}) = f(x_{r}) \implies x_{0} = x_{r}
    \end{equation*}
    Since \( f(x_{0}) = y_{0} = f(x_{r}) \), \( x_{0} = x_{r} \) follows by \aref{rule-of-inference:modus-ponens}{Modus Ponens}.
    Since \( x_{r} = x_{0} \) and \( x_{0} \in z_{0} \), it follows that \( x_{r} \in z_{0} \) as well. In other words,
    \begin{equation*}
        \begin{aligned}
            z_{0} \in S \implies x_{r} \in z_{0}&
        \end{aligned}
        \begin{aligned}
            &\text{for any \( z_{0} \)}
        \end{aligned}
    \end{equation*}
    Since \( z_{0} \) is arbitrary, we conclude by \aref{rule-of-inference:universal-generalization}{Universal Generalization} that,
    \begin{equation*}
        \forall z \, \big( z \in S \implies x_{r} \in z \big)
    \end{equation*}
    By \cref{dfn:intersection-of-sets}, it follows that \( x_{r} \in \Intersect S \).
    This means that,
    \begin{equation*}
        \begin{aligned}
            \orderedpair{x_{r}}{y_{0}} \in f&
        \end{aligned}
        \begin{aligned}
            &\text{for some \( x_{r} \in \Intersect S \)}
        \end{aligned}
    \end{equation*}
    By \cref{dfn:image-of-a-set-of-ordered-pairs}, necessarily \( y_{0} \in \imageSet{f}{\Intersect S} \).
    This entire chain of inferences shows that one can derive \( y_{0} \in \imageSet{f}{\Intersect S} \) from
    the initial assumption \( y_{0} \in \Intersect T \). Therefore, we can conclude by the \aref{rule-of-inference:deduction-theorem}{Deduction Theorem} that,
    \begin{equation*}
        \begin{aligned}
            y_{0} \in \Intersect T \implies y_{0} \in \imageSet{f}{\Intersect S}&
        \end{aligned}
        \begin{aligned}
            &\text{for any \( y_{0} \)}
        \end{aligned}
    \end{equation*}
    By \aref{rule-of-inference:universal-generalization}{Universal Generalization} and \cref{dfn:subset}, it follows that \( \Intersect T \subseteq \imageSet{f}{\Intersect S} \)
    for any injective function \( f \). Naturally, the equality follows from the fact that \( \imageSet{f}{\Intersect S} \subseteq \Intersect T\)
    is also true.
\end{proof}

\begin{cor}
    \label{thm:for-injective-functions-the-image-of-A-intersect-B-is-equal-to-the-image-of-A-intersect-the-image-of-B}
    If \( f \) is an injective function and \( A \) and \( B \) are sets, then,
    \begin{equation*}
        \imageSet{f}{A \intersect B} = \imageSet{f}{A} \intersect \imageSet{f}{B}
    \end{equation*}
\end{cor}

\begin{thm}
    \label{thm:for-injective-functions-the-difference-of-image-is-equal-to-the-image-of-difference}
    If \( f \) is an injective function and \( A \) and \( B \) are sets, then,
    \begin{equation*}
        \imageSet{f}{A - B} = \imageSet{f}{A} - \imageSet{f}{B}
    \end{equation*}
\end{thm}

\begin{proof}
    Let \( f \) be an injective function and \( A \) and \( B \) be arbitrary sets.
    It follows from \cref{thm:the-difference-of-image-is-a-subset-of-the-image-of-difference} that,
    \begin{equation*}
        \imageSet{f}{A} - \imageSet{f}{B} \subseteq \imageSet{f}{A - B}
    \end{equation*}
    To establish the equality, we have to prove the inclusion relation in the opposite direction.

    \vspace{0.5cm}
    \noindent
    Let \( y_{0} \) be an arbitrary object. If \( y_{0} \in \imageSet{f}{A-B} \), then,
    \begin{equation}
        \label{eq:y0-equals-fx0}
        \begin{aligned}
            \orderedpair{x_{0}}{y_{0}} \in f&
        \end{aligned}
        \begin{aligned}
            &\text{for some \( x_{0} \in A \land x_{0} \notin B \)}
        \end{aligned}
    \end{equation}
    Clearly, \( y_{0} \in \imageSet{f}{A} \). So the only task left is to show that
    \( y_{0} \notin \imageSet{f}{B} \), which requires one to prove that,
    \begin{equation}
        \label{eq:y0-does-not-belong-to-image-of-f-under-B}
        \forall x \, \big( x \in B \implies \orderedpair{x}{y_{0}} \notin f \big)
    \end{equation}
    Let \( \beta \in B \) be an arbitrary. Suppose that \( \orderedpair{\beta}{y_{0}} \in f \).
    Then \( f(\beta) = y_{0} \) by \cref{dfn:functional-notation}. But from \eqref{eq:y0-equals-fx0},
    we have \( f(x_{0}) = y_{0} \). Therefore, \( f(\beta) = f(x_{0}) \). Since \( f \) is injective,
    \begin{equation*}
        f(\beta) = f(x_{0}) \implies \beta = x_{0}
    \end{equation*}
    Thus, \( \beta = x_{0} \) by \aref{rule-of-inference:modus-ponens}{Modus Ponens}. But if \( \beta = x_{0} \), then \( \beta \notin B \),
    an obvious contradiction.
    Thus, we conclude by Reductio Ad Absurdum that the supposition \( \orderedpair{\beta}{y_{0}} \in f \)
    must be false. And by the \aref{rule-of-inference:deduction-theorem}{Deduction Theorem}, we infer,
    \begin{equation*}
        \beta \in B \implies \orderedpair{\beta}{y_{0}} \notin f
    \end{equation*}
    Since \( \beta \) is arbitrary, it follows from \aref{rule-of-inference:universal-generalization}{Universal Generalization} that,
    \begin{equation*}
        \forall x \, \big( x \in B \implies \orderedpair{x}{y_{0}} \notin f \big)
    \end{equation*}
    Thus, \( y_{0} \notin \imageSet{f}{B} \). Since \( y_{0} \in \imageSet{f}{A} \) and \( y_{0} \notin \imageSet{f}{B} \),
    \( y_{0} \in \imageSet{f}{A} - \imageSet{f}{B} \) by \cref{dfn:difference}.
    \begin{equation*}
        \begin{aligned}
            \therefore \, y_{0} \in \imageSet{f}{A-B} \implies y_{0} \in \imageSet{f}{A} - \imageSet{f}{B}&
        \end{aligned}
        \begin{aligned}
            &\text{for any \( y_{0} \)}
        \end{aligned}
    \end{equation*}
    By \aref{rule-of-inference:universal-generalization}{Universal Generalization} and \cref{dfn:subset}, it follows that \( \imageSet{f}{A-B} \subseteq \imageSet{f}{A} - \imageSet{f}{B} \).
    \( \imageSet{f}{A-B} = \imageSet{f}{A} - \imageSet{f}{B} \) follows from the fact that \( \imageSet{f}{A} - \imageSet{f}{B} \subseteq \imageSet{f}{A-B} \)
    is also true.
\end{proof}

\begin{thm}
    \label{thm:for-injective-functions-the-image-of-the-symdiff-is-the-symdiff-of-the-images}
    If \( f \) is an injective function and \( A \) and \( B \) are sets, then,
    \begin{equation*}
        \imageSet{f}{A \symdiff B} = \imageSet{f}{A} \symdiff \imageSet{f}{B}
    \end{equation*}
\end{thm}

\begin{proof}
    Let \( f \) be an injective function and \( A \) and \( B \) be arbitrary sets. By
    \cref{set-identities:symmetric-difference-as-the-union-of-differences}, we have,
    \begin{equation*}
        A \symdiff B = \mleft( A - B \mright) \union \mleft( B - A \mright)
    \end{equation*}
    Therefore,
    \begin{equation*}
        \imageSet{f}{A \symdiff B} = \imageSet{f}{\mleft( A - B \mright) \union \mleft( B - A \mright)}
    \end{equation*}
    And by \cref{thm:image-of-A-union-B-under-R-is-the-union-of-the-image-of-A-under-R-and-the-image-of-B-under-R},
    \begin{equation*}
        \imageSet{f}{A \symdiff B} = \imageSet{f}{A - B} \union \imageSet{f}{B - A}
    \end{equation*}
    Since \( f \) is injective, \cref{thm:for-injective-functions-the-difference-of-image-is-equal-to-the-image-of-difference}
    can be applied, so the equality becomes,
    \begin{equation*}
        \imageSet{f}{A \symdiff B} = \mleft( \imageSet{f}{A} - \imageSet{f}{B} \mright) \union \mleft( \imageSet{f}{B} - \imageSet{f}{A} \mright)
    \end{equation*}
    And by \cref{dfn:symmetric-difference}, we have \( \imageSet{f}{A \symdiff B} = \imageSet{f}{A} \symdiff \imageSet{f}{B} \)
\end{proof}

\begin{thm}
    \label{thm:the-union-of-a-compatible-system-of-functions-is-a-function}
    If \( F \) is a compatible system of functions, then \( \Union F \) is a function and,
    \begin{equation*}
        \dom \Union F = \Union \setbuilderBig{\dom f}{f \in F}
    \end{equation*}
\end{thm}

\begin{proof}
    Let \( F \) be a compatible system of functions. Thus, \( F \neq \emptyset \) by definition
    and every element of \( F \) is a function. To prove that \( \Union F \) is a function, we must
    show that it satisfies \cref{dfn:function-basic}. In other words, the following two statements must hold:
    \begin{gather}
        \forall x \, \mleft( x \in \Union F \implies
        \exists a \, \exists b \, \mleft( x = \orderedpair{a}{b} \mright) \mright) \label{eq:Union-F-is-a-set-of-ordered-pairs} \\
        \forall x \, \forall y \, \forall z \, \mleft( \Big[\orderedpair{x}{y} \in \Union F
        \land \orderedpair{x}{z} \in \Union F \Big] \implies y = z \mright) \label{eq:Union-F-has-unique-first-coordinate}
    \end{gather}
    Proving \eqref{eq:Union-F-is-a-set-of-ordered-pairs} is quite trivial. Let \( \alpha \in \Union F \). By \cref{dfn:union-of-sets},
    it follows that \( \alpha \in f \) for some \( f \in F \). Since \( F \) is a system of functions, it follows that \( f \) must
    be a function. And since functions are sets of ordered pairs, \( \alpha \) must therefore be an ordered pair.

    \vspace{0.3cm}
    \noindent
    To prove \eqref{eq:Union-F-has-unique-first-coordinate}, we let \orderedpair{x_{0}}{y_{0}} and \orderedpair{x_{0}}{z_{0}} be
    ordered pairs that belong to \( \Union F \). Again, by \cref{dfn:union-of-sets}, it follows that,
    \begin{equation*}
        \begin{aligned}
            \orderedpair{x_{0}}{y_{0}} \in f& \\
            \orderedpair{x_{0}}{z_{0}} \in g&
        \end{aligned}
        \begin{aligned}
            &\quad\text{for some \( f \in F \)}\\
            &\quad\text{for some \( g \in F \)}
        \end{aligned}
    \end{equation*}
    Since \( \Union F \) is a compatible system of functions, \( f \) and \( g \) must be compatible with each other.
    Thus, by \cref{thm:compatibility-of-two-functions}, we have,
    \begin{equation}
        \label{eq:f-and-g-have-the-same-values-where-their-domains-intersect}
        \forall x \, \Big( x \in \mleft( \dom f \intersect \dom g \mright) \implies f(x) = g(x) \Big)
    \end{equation}
    Since \( \orderedpair{x_{0}}{y_{0}} \in f \land \orderedpair{x_{0}}{z_{0}} \in g \), clearly \( x_{0} \in \dom f \land x_{0} \in \dom g \).
    By \cref{thm:intersection-in-terms-of-logical-and}, necessarily \( x_{0} \in \mleft( \dom f \intersect \dom g \mright) \).
    And by \eqref{eq:f-and-g-have-the-same-values-where-their-domains-intersect}, we infer \( f(x_{0}) = y_{0} = g(x_{0}) = z_{0} \).
    Thus, we have by the \aref{rule-of-inference:deduction-theorem}{Deduction Theorem},
    \begin{equation*}
        \Big( \orderedpair{x_{0}}{y_{0}} \in \Union F \land \orderedpair{x_{0}}{z_{0}} \in \Union F \Big) \implies \Big( y_{0} = z_{0} \Big)
    \end{equation*}
    Since \( x_{0} \), \( y_{0} \), and \( z_{0} \) are arbitrary, \eqref{eq:Union-F-has-unique-first-coordinate} follows readily from
    \aref{rule-of-inference:universal-generalization}{Universal Generalization}.

    \vspace{0.3cm}
    \noindent
    To prove that \( \dom \Union F = \Union \setbuilderbig{\dom f}{f \in F} \), we let \( x_{0} \) is an arbitrary object.
    Thus the following series of biconditional statements must hold:
    \begin{align*}
        x_{0} \in \dom \Union F
        &\iff
        \exists y \, \Big[ \orderedpair{x_{0}}{y} \in \Union F \Big] \\
        &\iff
        \exists y \, \Big[ \exists f \, \big( f \in F \land \orderedpair{x_{0}}{y} \in f \big) \Big] \\
        &\iff
        \exists f \, \Big[ f \in F \land \exists y \, \big( \orderedpair{x_{0}}{y} \in f \big) \Big] \\
        &\iff
        \exists f \, \Big[ f \in F \land x_{0} \in \dom f \Big] \\
        &\iff
        \exists f \, \Big[ f \in F \land \exists z \, \big( z = \dom f \land x \in z \big) \Big]\\
        &\iff
        \exists z \, \Big[ \exists f \, \big( f \in F \land z = \dom f \big) \land x \in z \Big] \\
        &\iff
        \exists z \, \Big[ z \in \setbuilderbig{\dom f}{f \in F} \land x \in z \Big] \\
        &\iff
        x_{0} \in \Union \setbuilderBig{\dom f}{f \in F}
    \end{align*}
    The equality follows immediately by \aref{rule-of-inference:universal-generalization}{Universal Generalization} and the \nameref{axm:axiom-of-extensionality}.
\end{proof}

\begin{thm}
    \phantomsection \label{thm:bijection-between-sets-implies-bijection-between-powersets}
    If a bijection exists between two sets \( A \) and \( B \), then a bijection also exists
    between \( \powerset{A} \) and \( \powerset{B} \).
\end{thm}

\begin{proof}
    Let \( A \) and \( B \) be arbitrary sets and \( \bijection{h}{A}{B} \) be a bijection between them.
    Let the set \( g \) be defined as follows:
    \begin{equation}
        \label{eq:define-the-mapping-g-between-powersets}
        g = \setbuilderBig{\orderedpair{x}{y}}{x \subseteq A \,\,{\land}\,\, y = \imageSet{h}{x} }
    \end{equation}
    Then we must prove that \( g \) is a bijection between \( \powerset{A} \) and \( \powerset{B} \), which entails
    proving the following claims:
    \begin{enumerate*}[label=(\roman*)]
        \item \( g \) is a \aref{dfn:function-basic}{function} with \( \dom g = \powerset{A} \) and \( \ran g \subseteq \powerset{B} \), \phantomsection \label{enum:g-is-a-function}
        \item \( g \) is \aref{dfn:injection}{injective}, \phantomsection \label{enum:g-is-injective}
        \item \( g \) is \aref{dfn:surjection}{surjective}. \phantomsection \label{enum:g-is-surjective}
    \end{enumerate*}

    \vspace{0.5cm}
    \noindent
    {\scshape Claim} \ref{enum:g-is-a-function}: \( g \) is a function with domain \( \powerset{A} \) and range \( \powerset{B} \)

    \vspace{0.1cm}
    \noindent
    To prove that \( g \) is a function, we must show that \( g \) satisfies \cref{dfn:function-basic}.
    From \eqref{eq:define-the-mapping-g-between-powersets}, it is obvious that \( g \) is a set of ordered pairs. That is,
    \begin{equation}
        \label{eq:g-is-a-set-of-ordered-pairs}
        \forall x \, \big[ x \in g \implies \exists a \, \exists b \, \mleft( x = \orderedpair{a}{b} \mright) \big]
    \end{equation}
    Now suppose that \( x \), \( y \), and \( z \) be arbitrary objects such that \( \orderedpair{x}{y} \in g \) and \( \orderedpair{x}{z} \in g \).
    By \eqref{eq:define-the-mapping-g-between-powersets}, it follows that \( y = \imageSet{h}{x} \) and \( z = \imageSet{h}{x} \).
    Clearly, \( y = z \).
    Thus, we have by the \aref{rule-of-inference:deduction-theorem}{Deduction Theorem} and \aref{rule-of-inference:universal-generalization}{Universal Generalization},
    \begin{equation}
        \label{eq:the-same-object-in-the-domain-of-g-is-associated-with-the-same-object-in-the-range-of-g}
        \forall x \, \forall y \, \forall z \, \Big[ \big( \orderedpair{x}{y} \in g \,\,{\land}\,\, \orderedpair{x}{z} \in g \big) \implies \mleft( y = z \mright) \Big]
    \end{equation}
    Applying \aref{rule-of-inference:conjunction-introduction}{Conjunction Introduction} to \eqref{eq:g-is-a-set-of-ordered-pairs}
    and \eqref{eq:the-same-object-in-the-domain-of-g-is-associated-with-the-same-object-in-the-range-of-g} then yields \cref{dfn:function-basic}. In other words,
    \( g \) is indeed a function. Furthermore, if \( \exists y \, \mleft[ \orderedpair{x}{y} \in g \mright] \), then \eqref{eq:define-the-mapping-g-between-powersets}
    requires that \( x \subseteq A \) and thus \( x \in \powerset{A} \). This means that \( \dom g \subseteq \powerset{A} \).
    Conversely, if \( x \) is an arbitrary subset of \( A \), then the image \( \imageSet{h}{x} \) must exist
    by \cref{thm:existence-of-the-image-of-a-set-of-ordered-pairs-under-another-set}. In other words, \( x \in \powerset{A} {\implies} x \in \dom g \), or
    \( \powerset{A} \subseteq \dom g \). Since \( \dom g \subseteq \powerset{A} \) and \( \powerset{A} \subseteq \dom g \) are both true, it follows that
    \( \dom g = \powerset{A}\). Next, let \( y_{0} \) be an arbitrary element of \( \ran g \). By \cref{dfn:range-of-a-set-of-ordered-pairs}, it follows
    that \( \orderedpair{x_{0}}{y_{0}} \in g \) for some \( x_{0} \). But according to \eqref{eq:define-the-mapping-g-between-powersets},
    \( \orderedpair{x_{0}}{y_{0}} \in g \) if and only if \( x_{0} \subseteq A \) and \( y_{0} = \imageSet{h}{x_{0}} \). Since \( \ran h = B \), it follows
    that \( y_{0} = \imageSet{h}{x_{0}} \subseteq B \). In other words, \( y_{0} \in \powerset{B} \).
    Thus, we have by the \aref{rule-of-inference:deduction-theorem}{Deduction Theorem} \( y_{0} \in \ran g {\implies} y_{0} \in \powerset{B} \) for any
    arbitrary \( y_{0} \). Finally, we conclude by \aref{rule-of-inference:universal-generalization}{Universal Generalization} that \( \ran g \subseteq \powerset{B}\).

    \vspace{0.3cm}
    \noindent
    {\scshape Claim} \ref{enum:g-is-injective}: \( g \) is injective

    \vspace{0.1cm}
    \noindent
    Let \( x, y \in \powerset{A}\) and \( z \in \powerset{B} \) be arbitrary sets such that,
    \begin{equation}
        \label{eq:gx-and-gy-equals-z}
        \orderedpair{x}{z} \in g \,\,{\land}\,\, \orderedpair{y}{z} \in g
    \end{equation}
    By \eqref{eq:define-the-mapping-g-between-powersets}, it follows that \( z = \imageSet{h}{x} \) and \( z = \imageSet{h}{y} \).
    Clearly, \( \imageSet{h}{x} = \imageSet{h}{y} \). Thus,
    \begin{equation}
        \label{eq:hx-equals-hy}
        \forall a \, \big( a \in \imageSet{h}{x} \iff a \in \imageSet{h}{y} \big)
    \end{equation}
    Let \( \alpha \in x \) be arbitrary. Since \( x \subseteq A \) and \( \dom h = A \), it follows that \( h(\alpha) \) is defined and that
    \( h(\alpha) \in \imageSet{h}{x} \). By \eqref{eq:hx-equals-hy}, it follows that \( h(\alpha) \in \imageSet{h}{y} \), and since \( h(\alpha) \in \imageSet{h}{y} \),
    \cref{dfn:image-of-a-set-of-ordered-pairs} guarantees that there exists some \( \beta \in y \) such that \( h(\alpha) = h(\beta) \).
    But \( h \) is a bijection between \( A \) and \( B \), so \cref{dfn:bijection} implies that \( \alpha = \beta \).
    Since \( \beta \in y \) and \( \alpha = \beta \), it follows that \( \alpha \in y \) must also be true. In other words,
    \( \alpha \in x {\implies} \alpha \in y \) for any arbitrary \( \alpha \). By \aref{rule-of-inference:universal-generalization}{Universal Generalization}, we have,
    \begin{equation*}
        \forall a \, \big( a \in x \implies a \in y \big)
    \end{equation*}
    The converse \( \forall a \, \mleft( a \in y {\implies} a \in x \mright) \) can be proven in the exact same way with the roles of \( x \) and \( y \) swapped.
    This shows that \( x = y \) and thus \( \orderedpair{x}{z} \in g \,\,{\land}\,\, \orderedpair{y}{z} \in g {\implies} x = y \) for any arbitrary \( x \) and \( y \).
    By \cref{dfn:injection}, it follows that \( g \) is an injective function.

    \vspace{0.3cm}
    \noindent
    {\scshape Claim} \ref{enum:g-is-surjective}: \( g \) is surjective

    \vspace{0.1cm}
    \noindent
    It has already been proven that \( \ran g \subseteq \powerset{B} \). Thus, it remains to show that \( \powerset{B} \subseteq \ran g \)
    to prove that \( g \) is surjective. Let \( y \in \powerset{B} \) be arbitrary.
    Since \( h \) is a bijection with \( \ran h = B \) and \( y \subseteq B \),
    it follows from \cref{thm:image-of-f-inverse-A-under-f-is-A-if-f-is-a-bijection-and-A-subset-of-range-f} that
    \( \imageSet{h}{\imageSet{h^{-1}}{y}} = y \).
    Furthermore, \( \imageSet{h^{-1}}{y} \subseteq \imageSet{h^{-1}}{B} \) by
    \cref{thm:A-subset-B-implies-image-of-A-under-R-subset-image-of-B-under-R}.
    Since \( h \) is a bijection, \cref{thm:the-inverse-of-a-bijection-is-also-a-bijection} guarantees that \( h^{-1} \) is also
    bijection where \( \dom h^{-1} = B \) and \( \ran h^{-1} = A \). Since \( \dom h^{-1} = B \), it follows that \( \imageSet{h^{-1}}{B} = A \)
    and thus \( \imageSet{h^{-1}}{y} \subseteq A \) and \( \imageSet{h^{-1}}{y} \in \powerset{A} \).
    By \eqref{eq:define-the-mapping-g-between-powersets},
    it follows that \( \orderedpair{\imageSet{h^{-1}}{y}}{y} \in g \) for any \( y \in \powerset{B} \).
    In other words, \( y \in \powerset{B} {\implies} y \in \ran g \), and by \aref{rule-of-inference:universal-generalization}{Universal Generalization},
    \( \powerset{B} \subseteq \ran g \) as needed.
\end{proof}

\newpage
\section{Indexed System of Sets}

\begin{dfn}[Indexed Sets]
    \label{dfn:indexed-system-of-sets}
    A set \( S \) is \keyword{an indexed system of sets}
    if and only if \surjection{S}{I}{R} for some set \( I \) and
    a system of sets \( R \). In such cases, we also say that \( R \)
    is \keyword{indexed} by \( S \). The notations \( \sequence{ S_{i} }{ i \in I } \)
    and \( \setbuilder{S_{i}}{i \in I } \) can be used to denote \( S \) and \( \ran S \)
    respectively.
\end{dfn}

\begin{dfn}[Set Exponentiation]
    \label{dfn:set-exponentiation}
    For any two sets \( A \) and \( B \), the set \( A^{B} \) is defined as,
    \begin{equation*}
        A^{B} = \setbuilderbig{f}{\function{f}{B}{A}}
    \end{equation*}
    That is, \( A^{B} \) is the set of all functions on \( B \) into \( A \).
\end{dfn}

\begin{dfn}[Union]
    \label{dfn:union-of-indexed-system-of-sets}
    If \surjection{S}{I}{R} is an indexed system of sets, then the \keyword{union} of \( S \) is defined
    as,
    \begin{equation*}
        \Union_{i \,\in\, I} S_{i}
        = \Union \ran S
        = \Union \setbuilderBig{S_{i}}{i \in I}
        = \setbuilderBig{x}{\exists i \, \mleft( i \in I \land x \in S_{i} \mright)}
    \end{equation*}
\end{dfn}

\begin{dfn}[Intersection]
    \label{dfn:intersection-of-indexed-system-of-sets}
    If \surjection{S}{I}{R} is an indexed system of sets with \( I \) and \( R \) being non-empty,
    then the \keyword{intersection} of \( S \) is defined as,
    \begin{equation*}
        \Intersect_{i \,\in\, I} S_{i}
        = \Intersect \ran S
        = \Intersect \setbuilderBig{S_{i}}{i \in I}
        = \setbuilderBig{x}{\forall i \, \big( i \in I \implies x \in S_{i} \big)}
    \end{equation*}
\end{dfn}

\begin{dfn}[Product]
    \label{dfn:product-of-indexed-system-of-sets}
    If \surjection{S}{I}{R} is an indexed system of sets, then the \keyword{product} of \( S \) is defined as,
    \begin{equation*}
        \prod S =
        \prod \sequence{S_{i}}{i \in I} =
        \prod_{i \,\in\, I} S_{i} = \setbuilderbigg{\function{f}{I}{\Union_{i \,\in\, I} S_{i}}}{\forall i \, \big( i \in I \implies f(i) \in S_{i} \big)}
    \end{equation*}
\end{dfn}

\vspace{0.1cm}
\noindent
We now prove some theorems involving indexed systems of sets.

\begin{thm}
    \label{thm:product-of-an-empty-set}
    \( \prod \emptyset = \mleft\{ \emptyset \mright\} \).
\end{thm}

\begin{proof}
    Since the product is only defined in the context of an indexed system of sets, we must treat \( \emptyset \)
    as an indexed system of sets. But there is only one possible indexed system of sets that is equal to \( \emptyset \):
    \begin{equation*}
        \surjection{S}{\emptyset}{\emptyset}
    \end{equation*}
    Thus, \( \prod \emptyset = \prod S \), which is the set,
    \begin{equation*}
        \setbuilderbigg{\function{f}{\emptyset}{\emptyset}}{\forall i \, \Big[ i \in \emptyset \implies
        \forall y \, \Big( \orderedpair{i}{y} \in f \implies \forall z \, \big[ \orderedpair{i}{z} \in S \implies y \in z \big] \Big)\Big]}
    \end{equation*}
    By \cref{thm:f-is-a-function-with-empty-domain-and-empty-range-if-and-only-if-it-is-the-empty-set}, the only set that can belong to
    \( \prod S \) is \( \emptyset \).
    \begin{equation*}
        \therefore \, \prod \emptyset = \mleft\{ \emptyset \mright\} \qedhere
    \end{equation*}
\end{proof}

\begin{thm}
    \label{thm:any-indexed-system-of-sets-with-empty-domain-must-have-empty-range-and-vice-versa}
    If \surjection{S}{I}{R} is an indexed system of sets, then \( I = \emptyset \iff R = \emptyset \).
\end{thm}

\begin{proof}
    Let \surjection{S}{I}{R} be an indexed system of sets. Suppose that \( I = \emptyset \). By
    \cref{dfn:surjection},
    \begin{equation}
        \label{eq:R-equals-ran-S}
        \forall y \, \Big[ y \in R \iff \exists x \, \big( x \in I \land \orderedpair{x}{y} \in S \big) \Big]
    \end{equation}
    If \( R \neq \emptyset \), then there must exist some \( y_{0} \in R \). But by \eqref{eq:R-equals-ran-S},
    this would mean that \( \orderedpair{x_{0}}{y_{0}} \) for some \( x_{0} \in I \), which contradicts the fact that
    \( I = \emptyset \). Thus, \( R = \emptyset \) by \aref{rule-of-inference:reductio-ad-absurdum}{Reductio ad Absurdum}.
    \begin{equation*}
        \therefore \, I = \emptyset \implies R = \emptyset
    \end{equation*}
    To prove the converse, suppose that \( R = \emptyset \). If \( I \neq \emptyset \), then
    there must exist some \( x_{0} \in I \). But \( I = \dom S \), so necessarily \( x_{0} \in \dom S \).
    By \cref{dfn:domain-of-a-set-of-ordered-pairs}, this means that there exists some \( y_{0} \in \ran S = R \)
    such that \( \orderedpair{x_{0}}{y_{0}} \in S \). This contradicts the initial supposition that \( R = \emptyset \).
    Again, by \aref{rule-of-inference:reductio-ad-absurdum}{Reductio ad Absurdum}, \( I \) must be empty.
    \begin{equation*}
        \therefore \, R = \emptyset \implies I = \emptyset \qedhere
    \end{equation*}
\end{proof}

\begin{thm}
    \label{thm:any-system-of-sets-can-be-indexed}
    If \( A \) is a system of sets, then there exists a function \( S \) such that
    \( A \) is indexed by \( S \).
\end{thm}

\begin{proof}
    Let \( S \) be the set defined as follows:
    \begin{equation*}
        S = \setbuilderbig{\orderedpair{x}{y}}{x \in A \land y = x}
    \end{equation*}
    We now prove that \( S \) is indeed a function. Let \( x_{0} \), \( y_{0} \), and \( z_{0} \)
    be arbitrary objects such that \( \orderedpair{x_{0}}{y_{0}} \in S \) and \( \orderedpair{x_{0}}{z_{0}} \in S \).
    By the definition of \( S \), it follows that \( y_{0} = x_{0} = z_{0} \). By \aref{rule-of-inference:universal-generalization}{Universal Generalization},
    \begin{equation*}
        \forall x \, \forall y \, \forall z \, \Big[ \orderedpair{x}{y} \in S \land \orderedpair{x}{z} \in S \implies y = z \Big]
    \end{equation*}
    By \cref{dfn:function-basic}, it follows that \( S \) is a function. It is trivial to prove that \( \dom f = A \) and \( \ran f = A \).
    Thus,
    \begin{equation*}
        \surjection{S}{A}{A}
    \end{equation*}
    It is evident that \( S \) is an indexed system of sets with \( I = A \) and \( R = A \).
\end{proof}

\begin{thm}
    \label{thm:existence-of-set-exponentiation}
    For any two sets \( A \) and \( B \), the set \( A^{B} \) exists.
\end{thm}

\begin{proof}
    Let the predicate \( \varphi \) be defined as,
    \begin{equation*}
        \varphi(f) \coloneq \function{f}{B}{A}
    \end{equation*}
    Let \( f \) be an arbitrary object such that \( \varphi(f) \) is satisfied. Thus,
    \begin{equation*}
        \function{f}{B}{A}
    \end{equation*}
    By \cref{thm:a-function-on-A-into-B-is-a-subset-of-the-cartesian-product-between-A-and-B},
    it follows that \( f \subseteq B \times A \), and thus \( f \in \powerset{B \times A} \).
    Since \( f \) is arbitrary, we infer by \aref{rule-of-inference:universal-generalization}{Universal Generalization},
    \begin{equation*}
        \forall f \, \Big[ \varphi(f) \implies f \in \powerset{B \times A} \Big]
    \end{equation*}
    Thus, the set \( A^{B} = \setbuilderbig{f}{\varphi(f)} = \setbuilderbig{f}{\function{f}{B}{A}} \) exists
    by \cref{thm:sufficient-condition-for-unrestricted-comprehension}.
\end{proof}

\begin{thm}
    \label{thm:the-exponentiation-of-the-empty-set}
    For any set \( A \),
    \begin{equation*}
        \emptyset^{A} =
        \begin{cases}
            \mleft\{ \emptyset \mright\}  & A = \emptyset \\
            \,\,\,\emptyset               & A \neq \emptyset
        \end{cases}
    \end{equation*}
\end{thm}

\begin{proof}
    If \( A = \emptyset \), then by \cref{dfn:set-exponentiation},
    \begin{equation*}
        \emptyset^{A} = \emptyset^\emptyset = \setbuilder{f}{\function{f}{\emptyset}{\emptyset}}
    \end{equation*}
    And by \cref{thm:f-is-a-function-with-empty-domain-and-empty-range-if-and-only-if-it-is-the-empty-set},
    \( \function{f}{\emptyset}{\emptyset} \iff f = \emptyset \). Thus,
    \begin{equation*}
        \emptyset^{\emptyset} = \mleft\{ \emptyset \mright\}
    \end{equation*}
    On the other hand, if \( A \neq \emptyset \), then \( \function{f}{A}{\emptyset} \) can never be satisfied
    because \( A \neq \emptyset \) implies that \( \ran f \neq \emptyset \). Thus, \( \emptyset^{A} = \emptyset \).
\end{proof}

\begin{thm}
    \phantomsection \label{thm:set-exponentiation-with-empty-exponent}
    For any set \( A \), \( A^{0} = A^{\emptyset} = \mleft\{ \emptyset \mright\} \).
\end{thm}

\begin{proof}
    Let \( A \) be an arbitrary set. It follows from \cref{thm:empty-set-universal-subset} that \( \emptyset \subseteq A \).
    Moreover, we have by \cref{thm:f-is-a-function-with-empty-domain-and-empty-range-if-and-only-if-it-is-the-empty-set},
    \( \function{\emptyset}{\emptyset}{\emptyset} \). In other words, \( \ran \emptyset \subseteq A \) and thus
    \( \function{\emptyset}{\emptyset}{A} \) by \cref{dfn:arrow-notation}, which in turn implies \( \emptyset \in A^{\emptyset} \) by \cref{dfn:set-exponentiation}.
    By the \aref{rule-of-inference:deduction-theorem}{Deduction Theorem} and \aref{rule-of-inference:universal-generalization}{Universal Generalization},
    we conclude that \( \forall x \, \mleft[ x = \emptyset {\implies} x \in A^{\emptyset} \mright] \).

    \vspace{0.3cm}
    \noindent
    Conversely, if \( f \in A^{\emptyset} \), then \( \function{f}{\emptyset}{A} \) by \cref{dfn:set-exponentiation}.
    Since \( \dom f = \emptyset \), necessarilly \( \ran f = \emptyset \). In other words, \( \function{f}{\emptyset}{\emptyset} \),
    and by \cref{thm:f-is-a-function-with-empty-domain-and-empty-range-if-and-only-if-it-is-the-empty-set}, necessarily \( f = \emptyset \).
    By the \aref{rule-of-inference:deduction-theorem}{Deduction Theorem}, we infer that \( f \in A^{\emptyset} {\implies} f = \emptyset \).
    Since \( f \) is arbitrary, it follows from \aref{rule-of-inference:universal-generalization}{Universal Generalization} that
    \( \forall x \, \mleft[ x \in A^{\emptyset} {\implies} x = \emptyset \mright] \).

    \vspace{0.3cm}
    \noindent
    Since \( \forall x \, \mleft[ x = \emptyset {\implies} x \in A^{\emptyset} \mright]  \) and
    \( \forall x \, \mleft[ x \in A^{\emptyset} {\implies} x = \emptyset \mright] \) are both true, we conclude by
    \aref{rule-of-inference:biconditional-introduction}{Biconditional Introduction} that,
    \begin{equation*}
        \forall x \, \big[ x \in A^{\emptyset} \iff x = \emptyset \big]
    \end{equation*}
    In other words, \( A^{\emptyset} \) is the singleton \( \mleft\{ \emptyset \mright\} \).
\end{proof}

\begin{thm}
    \label{thm:existence-of-the-union-of-an-indexed-system-of-sets}
    If \surjection{S}{I}{R} is an indexed system of sets, then the union of \( S \) exists.
\end{thm}

\begin{proof}
    Let \surjection{S}{I}{R} be an indexed system of sets. By \cref{dfn:union-of-indexed-system-of-sets}, the
    union of \( S \) is defined as,
    \begin{equation*}
        \Union_{i \,\in\, I} S_{i} = \Union \ran S = \Union R
    \end{equation*}
    Since \( R \) is an existing set, the \nameref{axm:axiom-of-union} guarantees the existence of \( \Union_{i \,\in\, I} S_{i} \).
\end{proof}

\begin{thm}
    \label{thm:existence-of-the-intersection-of-an-indexed-system-of-sets}
    If \surjection{S}{I}{R} is an indexed system of sets where \( I \neq \emptyset \) and \( R \neq \emptyset \),
    then the intersection of \( S \) exists.
\end{thm}

\begin{proof}
    Let the predicate \( \varphi \) be defined as,
    \begin{equation*}
        \varphi(x) \coloneq \forall i \, \Big[ i \in I \implies x \in S_{i} \Big]
    \end{equation*}
    Let \( x_{0} \) be an arbitrary object such that \( \varphi(x_{0}) \) is satisfied. Thus,
    \begin{equation}
        \label{eq:x0-belongs-to-all-Si}
        \forall i \, \Big[ i \in I \implies x_{0} \in S_{i} \Big]
    \end{equation}
    By the \nameref{axm:axiom-of-union}, \( \Union R \) exists. Since \( I \neq \emptyset \),
    it follows that there exists an \( i \) such that \( i \in I \). Since \( \ran S = R \),
    it follows that \( S_{i} \in R \). Since \( x_{0} \in S_{i} \) by \eqref{eq:x0-belongs-to-all-Si},
    and since \( S_{i} \in R \), it follows that \( x_{0} \in \Union R \). By the \aref{rule-of-inference:deduction-theorem}{Deduction Theorem},
    we conclude,
    \begin{equation*}
        \varphi(x_{0}) \implies x_{0} \in \Union R
    \end{equation*}
    Since \( x_{0} \) is arbitrary, \aref{rule-of-inference:universal-generalization}{Universal Generalization} yields,
    \begin{equation*}
        \forall x \, \Big[ \varphi(x) \implies x \in \Union R \Big]
    \end{equation*}
    By \cref{thm:sufficient-condition-for-unrestricted-comprehension}, the set
    \( \Intersect_{i \,\in\, I} S_{i} = \setbuilderbig{x}{\varphi(x)} = \setbuilderbig{x}{\forall i \, \mleft( i \in I \implies x \in S_{i} \mright)} \)
    exists.
\end{proof}

\begin{thm}
    \label{thm:existence-of-the-product-of-an-indexed-system-of-sets}
    If \surjection{S}{I}{R} is an indexed system of sets, then the product of \( S \) exists.
\end{thm}

\begin{proof}
    Let \surjection{S}{I}{R} be an indexed system of sets. By \cref{thm:existence-of-the-union-of-an-indexed-system-of-sets},
    \( \Union_{i \,\in\, I} S_{i} \) exists. Let the predicate \( \varphi \) be defined as,
    \begin{equation*}
        \varphi(f) \coloneq \function{f}{I}{\Union_{i \,\in\, I} S_{i}}
        \land
        \forall i \, \big( i \in I \implies f(i) \in S_{i} \big)
    \end{equation*}
    Let \( f \) be an arbitrary object such that \( \varphi(f) \) is satisfied. Thus, by the definition of \( \varphi \),
    \begin{equation*}
        \function{f}{I}{\Union_{i \,\in\, I} S_{i}} \land \forall i \, \big( i \in I \implies f(i) \in S_{i} \big)
    \end{equation*}
    Since \function{f}{I}{\Union_{i \,\in\, I} S_{i}},
    it follows from \cref{thm:a-function-on-A-into-B-is-a-subset-of-the-cartesian-product-between-A-and-B}
    that \( f \subseteq I \times \Union_{i \,\in\, I} S_{i} \). In other words, \( f \in \powerset{I \times \Union_{i \,\in\, I} S_{i}} \).
    Thus,
    \begin{equation*}
        \begin{aligned}
            \varphi(f) \implies f \in \powerset{I \times \Union_{i \,\in\, I} S_{i}}&
        \end{aligned}
        \begin{aligned}
            &\text{for any \( f \)}
        \end{aligned}
    \end{equation*}
    By \cref{thm:sufficient-condition-for-unrestricted-comprehension}, the set \( \prod S = \setbuilder{f}{\varphi(f)} \) exists.
\end{proof}

\begin{thm}
    \label{thm:product-of-an-indexed-singleton}
    If \( A \) is a set and \surjection{S}{I}{\mleft\{ A \mright\}} for some non-empty index set \( I \), then
    \begin{equation*}
        \prod_{i \,\in\, I} S_{i} = A^{I}
    \end{equation*}
\end{thm}

\begin{proof}
    Let \( A \) be a set. Let \( S \) be a function that indexes the singleton \( \mleft\{ A \mright\} \). That is,
    \surjection{S}{I}{\mleft\{ A \mright\}} for some non-empty index set \( I \).
    By \cref{dfn:product-of-indexed-system-of-sets,thm:existence-of-the-product-of-an-indexed-system-of-sets},
    the product of \( S \) exists and is given by:
    \begin{equation}
        \label{eq:express-the-product-of-S-using-set-builder-notation}
        \prod_{i \,\in\, I} S_{i} = \setbuilderBig{\function{f}{I}{\Union_{i \,\in\, I} S_{i} }}{\forall i \, \big( i \in I \implies f(i) \in S_{i} \big)}
    \end{equation}
    But since \( \ran S = \mleft\{ A \mright\} \), it follows that,
    \begin{equation*}
        \forall i \, \big( i \in I \implies S_{i} = A \big)
    \end{equation*}
    Also, \( \Union_{i \,\in\, I} S_{i} = \Union \mleft\{ A \mright\} = A \) by \cref{thm:union-of-a-singleton}.
    Thus, \eqref{eq:express-the-product-of-S-using-set-builder-notation} becomes,
    \begin{equation}
        \label{eq:the-product-of-S-in-set-builder-notation-final-form}
        \prod_{i \,\in\, I} S_{i} = \setbuilderBig{\function{f}{I}{A}}{\forall i \, \big( i \in I \implies f(i) \in A \big)}
    \end{equation}
    Let \( f \) be an arbitrary object. If \( f \in \prod_{i \,\in\, I} S_{i} \), then by \eqref{eq:the-product-of-S-in-set-builder-notation-final-form},
    we have,
    \begin{equation*}
        \function{f}{I}{A} \land \forall i \, \big( i \in I \implies f(i) \in A \big)
    \end{equation*}
    Since \function{f}{I}{A}, clearly \( f \in A^{I} \) by \cref{dfn:set-exponentiation}.
    By the \aref{rule-of-inference:deduction-theorem}{Deduction Theorem},
    \begin{equation}
        \label{eq:f-in-product-implies-f-in-AI}
        f \in \prod_{i \,\in\, I} S_{i} \implies f \in A^{I}
    \end{equation}
    Conversely, if \( f \in A^{I} \), then \function{f}{I}{A} by \cref{dfn:set-exponentiation}.
    Thus, by \cref{thm:full-definition-of-the-arrow-notation},
    \begin{equation}
        \label{eq:every-i-in-I-has-a-unique-value}
        \forall i \, \Big[ i \in I \implies \exists! y \, \big( y \in A \land \orderedpair{i}{y} \in f \big) \Big]
    \end{equation}
    Let \( i \in I \) be arbitrary. Thus, by \eqref{eq:every-i-in-I-has-a-unique-value}, \( \orderedpair{i}{y_{0}} \in f \)
    for some unique \( y_{0} \in A \). Since \( y_{0} = f(i) \) by \cref{dfn:functional-notation}, necessarily \( f(i) \in A \).
    Thus,
    \begin{equation*}
        i \in I \implies f(i) \in A
    \end{equation*}
    Since \( i \) is arbitrary, \aref{rule-of-inference:universal-generalization}{Universal Generalization} yields,
    \begin{equation*}
        \forall i \, \Big[ i \in I \implies f(i) \in A \Big]
    \end{equation*}
    Since \function{f}{I}{A} and \( \forall i \, \mleft( i \in I \implies f(i) \in A \mright) \) are both true,
    \( f \in \prod_{i \,\in\, S_{i}} \) follows by \eqref{eq:the-product-of-S-in-set-builder-notation-final-form}.
    Again, by the \aref{rule-of-inference:deduction-theorem}{Deduction Theorem},
    \begin{equation}
        \label{eq:f-in-AI-implies-f-in-product}
        f \in A^{I} \implies f \in \prod_{i \,\in\, I} S_{i}
    \end{equation}
    Combining \eqref{eq:f-in-product-implies-f-in-AI} and \eqref{eq:f-in-AI-implies-f-in-product},
    we have for any arbitrary \( f \),
    \begin{equation*}
            f \in A^{I} \iff f \in \prod_{i \,\in\, I} S_{i}
    \end{equation*}
    The equality follows immediately from \aref{rule-of-inference:universal-generalization}{Universal Generalization} and the \nameref{axm:axiom-of-extensionality}.
\end{proof}
